\documentclass[11pt,a4paper]{article}
\usepackage[utf8]{inputenc}
\usepackage[T1]{fontenc}
\usepackage[french]{babel}
\usepackage{amsmath,amssymb,amsfonts,amsthm,mathrsfs,mathtools}
\usepackage[top=2cm,bottom=2cm,left=2cm,right=2cm]{geometry}
\usepackage[dvipsnames]{xcolor}
\usepackage{tcolorbox}
\tcbuselibrary{skins,breakable}
\usepackage{tikz}
\usetikzlibrary{calc,shapes,arrows.meta}
\usepackage{pgfplots}
\pgfplotsset{compat=1.18}
\usepackage{enumitem}
\usepackage{fancyhdr}
\usepackage{lastpage}
\usepackage{hyperref}

% ----------------------------------------------------
% Couleurs officielles du projet (Charte LLG / MPSI)
% ----------------------------------------------------
\definecolor{llgblue}{RGB}{20, 50, 115}
\definecolor{llglight}{RGB}{246, 249, 255}

\definecolor{emeraldframe}{RGB}{18, 110, 75}
\definecolor{emeraldback}{RGB}{242, 250, 246}
\definecolor{emeraldtitle}{RGB}{12, 80, 52}

\definecolor{prepaframe}{RGB}{30, 70, 135}
\definecolor{prepaback}{RGB}{246, 249, 255}
\definecolor{prepatitle}{RGB}{20, 50, 105}

\definecolor{alertframe}{RGB}{185, 75, 20}
\definecolor{alertback}{RGB}{255, 249, 245}
\definecolor{alerttitle}{RGB}{145, 55, 12}

\definecolor{bonusframe}{RGB}{120, 40, 140}
\definecolor{bonusback}{RGB}{252, 247, 255}
\definecolor{bonustitle}{RGB}{90, 25, 110}

% Configuration des hyperliens
\hypersetup{
    colorlinks=true,
    linkcolor=llgblue,
    urlcolor=llgblue,
    citecolor=llgblue
}

% En-têtes et pieds de page
\pagestyle{fancy}
\fancyhf{}
\renewcommand{\headrulewidth}{0.5pt}
\renewcommand{\footrulewidth}{0.4pt}
\fancyhead[L]{\small\textbf{Lycée Louis-le-Grand} -- 1\textsuperscript{ère} Spécialité Mathématiques (1E3)}
\fancyhead[C]{\small\textbf{Travaux Dirigés : La Fonction Exponentielle}}
\fancyhead[R]{\small Février 2022}
\fancyfoot[L]{\footnotesize LLG 2021/2022 -- Feuille d'exercices n°5}
\fancyfoot[C]{\footnotesize Énoncés \& Corrigés Exhaustifs d'Excellence}
\fancyfoot[R]{\footnotesize Page \thepage\ / \pageref{LastPage}}

% ----------------------------------------------------
% Boîtes tcolorbox personnalisées
% ----------------------------------------------------
\newtcolorbox{theoremebox}[1][]{
  colback=llglight,
  colframe=llgblue,
  fonttitle=\bfseries\color{white},
  coltitle=white,
  colbacktitle=llgblue,
  title={Théorème fondamental & Définition canonique #1},
  arc=2mm,
  boxrule=0.8pt,
  titlerule=0.5pt,
  breakable,
  before skip=0.4cm,
  after skip=0.5cm
}

\newtcolorbox{methodebox}[1][]{
  colback=prepaback,
  colframe=prepaframe,
  fonttitle=\bfseries\color{white},
  coltitle=white,
  colbacktitle=prepatitle,
  title={Méthodologie & Réflexe de Rédaction #1},
  arc=2mm,
  boxrule=0.8pt,
  titlerule=0.5pt,
  breakable,
  before skip=0.4cm,
  after skip=0.5cm
}

\newtcolorbox{correctionbox}[1][]{
  colback=emeraldback,
  colframe=emeraldframe,
  fonttitle=\bfseries\color{white},
  coltitle=white,
  colbacktitle=emeraldtitle,
  title={Corrigé détaillé et rédigé pas à pas #1},
  arc=2mm,
  boxrule=0.8pt,
  titlerule=0.5pt,
  breakable,
  before skip=0.4cm,
  after skip=0.5cm
}

\newtcolorbox{remarquebox}[1][]{
  colback=prepaback,
  colframe=prepaframe,
  fonttitle=\bfseries\small\color{white},
  coltitle=white,
  colbacktitle=prepatitle,
  title={Prolongement & Regard Prépa / MPSI #1},
  arc=1.5mm,
  boxrule=0.6pt,
  titlerule=0.4pt,
  breakable,
  before skip=0.25cm,
  after skip=0.35cm
}

\newtcolorbox{alertebox}[1][]{
  colback=alertback,
  colframe=alertframe,
  fonttitle=\bfseries\small\color{white},
  coltitle=white,
  colbacktitle=alerttitle,
  title={Point de vigilance méthodologique & Pièges fréquents #1},
  arc=1.5mm,
  boxrule=0.6pt,
  titlerule=0.4pt,
  breakable,
  before skip=0.25cm,
  after skip=0.35cm
}

% Commandes mathématiques usuelles
\newcommand{\N}{\mathbb{N}}
\newcommand{\Z}{\mathbb{Z}}
\newcommand{\Q}{\mathbb{Q}}
\newcommand{\R}{\mathbb{R}}
\newcommand{\C}{\mathbb{C}}
\newcommand{\sh}{\operatorname{sh}}
\newcommand{\ch}{\operatorname{ch}}
\renewcommand{\th}{\operatorname{th}}

\title{\Huge \textbf{La Fonction Exponentielle}\\[0.4cm]
\Large Recueil intégral des Travaux Dirigés d'Approfondissement \& Corrigés Exhaustifs}
\author{\textbf{Classe de Première Spécialité Mathématiques (1E3)}\\
\textit{Lycée Louis-le-Grand -- Parcours d'Excellence vers les CPGE Scientifiques (MPSI/PCSI)}}
\date{Année scolaire 2021--2022 -- Février 2022}

\begin{document}

\maketitle

\vspace{-0.3cm}
\begin{abstract}
Ce document constitue la référence exhaustive du chapitre consacré à \textbf{la fonction exponentielle} pour la classe de Première 1E3 du Lycée Louis-le-Grand (année scolaire 2021--2022). Conçu pour consolider les fondations de l'analyse réelle et ouvrir sur les exigences des classes préparatoires scientifiques (MPSI/PCSI), il comprend :
\begin{enumerate}[label=\textbf{\Roman*.}]
    \item Un cours de synthèse approfondi explicitant la construction axiomatique (problème de Cauchy $y'=y, y(0)=1$), l'unicité et le morphisme fondamental $(\R,+) \to (\R^*_+,\times)$, les théorèmes de convexité et l'échelle des puissances ;
    \item L'énoncé complet et le \textbf{corrigé intégral pas à pas des 12 exercices officiels de travaux dirigés}, articulés en quatre sections thématiques : calculs algébriques, dérivation et variations, tangentes et familles de courbes, calculs de limites et levée d'indéterminations par croissances comparées ;
    \item Des encarts d'enrichissement méthodologique, des figures géométriques TikZ de haute précision et des prolongements vers la trigonométrie hyperbolique et les développements limités.
\end{enumerate}
\end{abstract}

\vspace{0.3cm}
\hrule
\vspace{0.3cm}
\tableofcontents
\vspace{0.5cm}
\hrule
\vspace{0.8cm}
\section{Synthèse théorique : Fondements axiomatiques et propriétés analytiques}

La fonction exponentielle constitue l'une des pierres angulaires de l'analyse mathématique moderne. Au Lycée Louis-le-Grand, l'étude de l'exponentielle ne se limite pas à une simple manipulation calculatoire ; elle est introduite par le biais de son équation différentielle constitutive (problème de Cauchy), mettant en lumière les liens intimes entre dérivation, structure algébrique de groupe et géométrie analytique.

\subsection{Définition différentielle et unicité par fonction auxiliaire}

\begin{theoremebox}[Théorème d'existence et d'unicité (Euler-Cauchy)]
Il existe une unique fonction $f$, dérivable sur $\R$, solution du problème différentiel :
\[
\begin{cases}
f'(x) = f(x) \quad \text{pour tout } x \in \R, \\
f(0) = 1.
\end{cases}
\]
Cette unique fonction est appelée \textbf{fonction exponentielle}, notée $\exp$, et l'on note pour tout réel $x$ :
\[
\exp(x) = e^x \quad \text{où } e = \exp(1) \approx 2{,}71828\,18284\dots
\]
\end{theoremebox}

\begin{proof}[\textbf{Démonstration intégrale de l'unicité et de la non-annulation}]
Soit $f$ une fonction dérivable sur $\R$ telle que $f' = f$ et $f(0) = 1$.

\textbf{Étape 1 : Non-annulation de $f$ sur $\R$.}\\
Considérons la fonction auxiliaire $\varphi : \R \to \R$ définie pour tout $x \in \R$ par :
\[
\varphi(x) = f(x) \times f(-x).
\]
Par composition et produit de fonctions dérivables sur $\R$, $\varphi$ est dérivable sur $\R$. Pour tout $x \in \R$, appliquons la règle de dérivation d'un produit $(uv)' = u'v + uv'$ combinée à la dérivation de $x \mapsto f(-x)$ de dérivée $-f'(-x)$ :
\[
\varphi'(x) = f'(x) f(-x) + f(x) \times \bigl(-f'(-x)\bigr) = f'(x) f(-x) - f(x) f'(-x).
\]
Comme $f$ satisfait l'équation différentielle $f' = f$, on a $f'(x) = f(x)$ et $f'(-x) = f(-x)$. Il s'ensuit que :
\[
\varphi'(x) = f(x) f(-x) - f(x) f(-x) = 0.
\]
La dérivée de $\varphi$ étant identiquement nulle sur l'intervalle $\R$, la fonction $\varphi$ est \textbf{constante} sur $\R$. Sa valeur en tout point coïncide avec sa valeur en $x = 0$ :
\[
\varphi(0) = f(0) \times f(-0) = 1 \times 1 = 1.
\]
Par conséquent, pour tout $x \in \R$ :
\[
f(x) f(-x) = 1.
\]
Cette relation fondamentale prouve que pour tout $x \in \R$, $f(x) \neq 0$ (car si un réel $x_0$ annulait $f$, le produit $f(x_0)f(-x_0)$ s'annulerait, contredisant $1 \neq 0$). De plus, elle fournit immédiatement l'identité de l'inverse :
\[
f(-x) = \frac{1}{f(x)}.
\]

\textbf{Étape 2 : Unicité de la solution.}\\
Supposons qu'il existe une seconde fonction $g$ dérivable sur $\R$ vérifiant $g' = g$ et $g(0) = 1$. Puisque $f$ ne s'annule jamais sur $\R$, nous pouvons considérer la fonction quotient $h : \R \to \R$ définie par :
\[
h(x) = \frac{g(x)}{f(x)}.
\]
$h$ est dérivable sur $\R$ comme quotient de fonctions dérivables dont le dénominateur ne s'annule pas. Calculons sa dérivée :
\[
h'(x) = \frac{g'(x) f(x) - g(x) f'(x)}{\bigl(f(x)\bigr)^2}.
\]
En utilisant $g' = g$ et $f' = f$, le numérateur devient :
\[
g(x) f(x) - g(x) f(x) = 0.
\]
Ainsi $h'(x) = 0$ pour tout $x \in \R$. La fonction $h$ est donc constante sur $\R$, et sa valeur vaut :
\[
h(0) = \frac{g(0)}{f(0)} = \frac{1}{1} = 1.
\]
D'où, pour tout $x \in \R$, $h(x) = 1 \iff g(x) = f(x)$. L'unicité est rigoureusement démontrée.
\end{proof}

\subsection{Morphisme algébrique et règles opératoires}

\begin{theoremebox}[Équation fonctionnelle fondamentale]
Pour tous réels $a$ et $b$ :
\[
e^{a+b} = e^a \times e^b.
\]
En langage de théorie des groupes, la fonction exponentielle réalise un \textbf{isomorphisme de groupes} du groupe additif $(\R, +)$ vers le groupe multiplicatif $(\R^*_+, \times)$.
\end{theoremebox}

\begin{proof}
Fixons un réel $b \in \R$. Considérons la fonction auxiliaire $\psi : \R \to \R$ définie par :
\[
\psi(x) = \frac{f(x+b)}{f(b)} = \frac{e^{x+b}}{e^b}.
\]
Cette fonction est bien définie car $e^b \neq 0$, et dérivable sur $\R$. Sa dérivée vaut :
\[
\psi'(x) = \frac{1}{e^b} \times (e^{x+b})' = \frac{1}{e^b} \times 1 \cdot e^{x+b} = \psi(x).
\]
De plus, en $x = 0$ :
\[
\psi(0) = \frac{e^{0+b}}{e^b} = \frac{e^b}{e^b} = 1.
\]
La fonction $\psi$ vérifie donc l'équation différentielle $\psi' = \psi$ et la condition initiale $\psi(0) = 1$. Par unicité de la solution du problème de Cauchy démontrée précédemment, $\psi$ coïncide nécessairement avec la fonction exponentielle $f$ :
\[
\forall x \in \R, \quad \psi(x) = e^x \iff \frac{e^{x+b}}{e^b} = e^x \iff e^{x+b} = e^x e^b.
\]
En posant $x = a$, on obtient bien $e^{a+b} = e^a e^b$.
\end{proof}

\begin{methodebox}[Formulaire de calcul algébrique fondamental]
Pour tous réels $x, y$ et pour tout entier relatif $n \in \Z$ :
\begin{enumerate}[label=\textbf{\arabic*.}]
    \item $e^0 = 1$ \quad et \quad $e^1 = e \approx 2{,}718$.
    \item $e^{-x} = \dfrac{1}{e^x}$ \quad et \quad $e^{x-y} = \dfrac{e^x}{e^y}$.
    \item $(e^x)^n = e^{nx}$ \quad (généralisable à tout exposant rationnel ou réel : $(e^x)^r = e^{rx}$).
    \item $\sqrt{e^x} = e^{\frac{x}{2}}$ \quad (car $(\sqrt{e^x})^2 = e^x$ et $\sqrt{e^x} > 0$).
\end{enumerate}
\end{methodebox}

\subsection{Propriétés analytiques : Stricte positivité, croissance et convexité}

\begin{theoremebox}[Stricte positivité et stricte monotonie]
\begin{enumerate}
    \item Pour tout réel $x \in \R$, on a :
    \[
    e^x > 0.
    \]
    \item La fonction exponentielle est strictement croissante sur $\R$.
    \item Pour tous réels $a$ et $b$ :
    \[
    e^a = e^b \iff a = b, \qquad e^a < e^b \iff a < b.
    \]
\end{enumerate}
\end{theoremebox}

\begin{proof}
\begin{enumerate}
    \item Pour tout $x \in \R$, écrivons $x = \frac{x}{2} + \frac{x}{2}$. Par le morphisme fondamental :
    \[
    e^x = e^{\frac{x}{2} + \frac{x}{2}} = \left(e^{\frac{x}{2}}\right)^2.
    \]
    Un carré de réel est toujours positif ou nul. Or, nous avons établi à l'étape 1 que l'exponentielle ne s'annule jamais sur $\R$. Par conséquent, $e^{x/2} \neq 0$, ce qui implique que son carré est \textbf{strictement positif} : $e^x > 0$.
    \item La dérivée de la fonction exponentielle est elle-même : $(\exp)'(x) = e^x$. Puisque $e^x > 0$ pour tout $x \in \R$, la dérivée est strictement positive sur l'intervalle $\R$. Ainsi, la fonction exponentielle est \textbf{strictement croissante} sur $\R$.
    \item La stricte croissance garantit l'injectivité de la fonction, d'où les équivalences d'égalité et d'inégalité stricte.
\end{enumerate}
\end{proof}

\begin{theoremebox}[Convexité et inégalité de la tangente à l'origine]
La fonction exponentielle est deux fois dérivable sur $\R$ et sa dérivée seconde vérifie :
\[
(\exp)''(x) = (\exp)'(x) = e^x > 0.
\]
La fonction exponentielle est donc \textbf{strictement convexe} sur $\R$. Sa courbe représentative $\mathcal{C}_{\exp}$ est située \textbf{strictement au-dessus} de toutes ses tangentes (sauf au point de contact).

En particulier, la tangente au point d'abscisse $x_0 = 0$ a pour équation :
\[
y = \exp'(0)(x - 0) + \exp(0) = 1 \cdot x + 1 = x + 1.
\]
On en déduit l'\textbf{inégalité de convexité fondamentale} :
\[
\forall x \in \R, \quad e^x \ge x + 1,
\]
avec égalité si et seulement si $x = 0$. En particulier, pour tout $x \in \R$, $e^x > x$.
\end{theoremebox}

\subsection{Dérivation des fonctions composées du type $\exp(u)$}

\begin{theoremebox}[Théorème de composition]
Soit $u$ une fonction dérivable sur un intervalle $I \subset \R$.
La fonction $f = \exp \circ u : x \mapsto e^{u(x)}$ est dérivable sur $I$, et sa dérivée est donnée par :
\[
\forall x \in I, \quad \bigl(e^{u(x)}\bigr)' = u'(x) \cdot e^{u(x)}.
\]
Puisque $e^{u(x)} > 0$ pour tout $x \in I$, le signe de la dérivée $\bigl(e^{u(x)}\bigr)'$ est \textbf{strictement identique au signe de la dérivée $u'(x)$}.
\end{theoremebox}

\subsection{Théorèmes de croissances comparées}

\begin{theoremebox}[L'échelle des puissances et croissances comparées]
Au voisinage de l'infini, la fonction exponentielle l'emporte sur toute fonction puissance :
\begin{enumerate}
    \item Limites fondamentales :
    \[
    \lim_{x \to +\infty} e^x = +\infty \qquad \text{et} \qquad \lim_{x \to -\infty} e^x = 0.
    \]
    \item Croissances comparées en $+\infty$ : Pour tout entier $n \ge 1$ :
    \[
    \lim_{x \to +\infty} \frac{e^x}{x^n} = +\infty \qquad \text{et} \qquad \lim_{x \to +\infty} x^n e^{-x} = 0.
    \]
    \item Croissances comparées en $-\infty$ : Pour tout entier $n \ge 1$ :
    \[
    \lim_{x \to -\infty} x^n e^x = 0.
    \]
    \item Taux d'accroissement en $0$ :
    \[
    \lim_{x \to 0} \frac{e^x - 1}{x} = 1.
    \]
\end{enumerate}
\end{theoremebox}

\begin{remarquebox}[Ouverture vers les CPGE : Trigonométrie hyperbolique et séries entières]
En classe préparatoire (MPSI/PCSI), la fonction exponentielle est également définie par la somme d'une série entière convergente sur $\C$ :
\[
e^z = \sum_{n=0}^{+\infty} \frac{z^n}{n!} = 1 + z + \frac{z^2}{2!} + \frac{z^3}{3!} + \dots
\]
Cette approche relie l'analyse réelle à l'analyse complexe via la célèbre formule d'Euler $e^{ix} = \cos x + i \sin x$.
Elle donne naissance aux fonctions \textbf{cosinus hyperbolique} et \textbf{sinus hyperbolique} :
\[
\ch x = \frac{e^x + e^{-x}}{2}, \qquad \sh x = \frac{e^x - e^{-x}}{2}, \qquad \th x = \frac{\sh x}{\ch x} = \frac{e^x - e^{-x}}{e^x + e^{-x}},
\]
qui satisfont l'identité fondamentale $\ch^2 x - \sh^2 x = 1$, au cœur de la géométrie de Lorentz et de la relativité restreinte.
\end{remarquebox}

\vspace{0.8cm}
\hrule
\vspace{0.8cm}
\section{Section 1 : Calculs algébriques fondamentaux}

Cette première section est consacrée à la maîtrise absolue des règles opératoires sur l'exponentielle. Au niveau Première Spécialité et dans l'optique des classes préparatoires, la virtuosité technique dans la manipulation des exposants est indispensable pour aborder sereinement les études de fonctions, les résolutions d'équations et le calcul différentiel.

% =========================================================================
% EXERCICE 1
% =========================================================================
\subsection*{Exercice 1 : Simplification d'expressions exponentielles}
\addcontentsline{toc}{subsection}{Exercice 1 : Simplification d'expressions exponentielles}

Pour tout réel $x \in \R$, simplifier au maximum et mettre sous la forme la plus épurée les six expressions suivantes :
\begin{enumerate}[label=\textbf{\arabic*.}]
    \item $A_1(x) = (e^x)^3 e^{-2x}$
    \item $A_2(x) = \dfrac{e^x + e^{-x}}{e^{-x}}$
    \item $A_3(x) = \dfrac{e^{2x+1}}{e^{-2x}}$
    \item $A_4(x) = \sqrt{e^{2x}}$
    \item $A_5(x) = \sqrt{e^{-4x}}$
    \item $A_6(x) = \left(\dfrac{e^x + e^{-x}}{2}\right)^2 - \left(\dfrac{e^x - e^{-x}}{2}\right)^2$
\end{enumerate}

\begin{correctionbox}[Exercice 1]
Détaillons pas à pas chaque simplification en explicitant systématiquement les propriétés algébriques mises en jeu.

\begin{enumerate}[label=\textbf{\arabic*.}]
    \item \textbf{Simplification de $A_1(x) = (e^x)^3 e^{-2x}$ :}
    \begin{itemize}
        \item En vertu de la propriété des puissances d'exponentielles $(e^a)^n = e^{na}$, on a :
        \[
        (e^x)^3 = e^{3x}.
        \]
        \item En utilisant ensuite la règle d'addition des puissances pour un produit de même base $e^u e^v = e^{u+v}$ :
        \[
        A_1(x) = e^{3x} \times e^{-2x} = e^{3x + (-2x)} = e^{3x - 2x} = e^x.
        \]
    \end{itemize}
    \textbf{Conclusion :} $\boxed{A_1(x) = e^x}$.

    \item \textbf{Simplification de $A_2(x) = \dfrac{e^x + e^{-x}}{e^{-x}}$ :}
    \begin{itemize}
        \item \textbf{Méthode 1 (Décomposition du numérateur terme à terme) :}
        Le dénominateur $e^{-x}$ est strictement positif pour tout $x \in \R$, donc l'expression est définie sur $\R$. On sépare la fraction :
        \[
        A_2(x) = \frac{e^x}{e^{-x}} + \frac{e^{-x}}{e^{-x}}.
        \]
        D'une part, $\dfrac{e^{-x}}{e^{-x}} = 1$. D'autre part, en appliquant le quotient de puissances $\dfrac{e^u}{e^v} = e^{u-v}$ :
        \[
        \frac{e^x}{e^{-x}} = e^{x - (-x)} = e^{x+x} = e^{2x}.
        \]
        D'où :
        \[
        A_2(x) = e^{2x} + 1.
        \]
        \item \textbf{Méthode 2 (Multiplication par la quantité conjuguée $e^x$) :}
        En multipliant le numérateur et le dénominateur par $e^x$ (ce qui est légitime car $e^x > 0$) :
        \[
        A_2(x) = \frac{(e^x + e^{-x})e^x}{e^{-x}e^x} = \frac{e^x e^x + e^{-x}e^x}{e^{-x+x}} = \frac{e^{2x} + e^0}{e^0} = \frac{e^{2x} + 1}{1} = e^{2x} + 1.
        \]
    \end{itemize}
    \textbf{Conclusion :} $\boxed{A_2(x) = e^{2x} + 1}$.

    \item \textbf{Simplification de $A_3(x) = \dfrac{e^{2x+1}}{e^{-2x}}$ :}
    \begin{itemize}
        \item Par la règle du quotient de deux exponentielles $\dfrac{e^a}{e^b} = e^{a-b}$ :
        \[
        A_3(x) = e^{(2x+1) - (-2x)} = e^{2x + 1 + 2x} = e^{4x+1}.
        \]
        \item On peut également écrire cette expression sous la forme $e^{4x+1} = e^1 \times e^{4x} = e (e^x)^4$.
    \end{itemize}
    \textbf{Conclusion :} $\boxed{A_3(x) = e^{4x+1} = e \cdot e^{4x}}$.

    \item \textbf{Simplification de $A_4(x) = \sqrt{e^{2x}}$ :}
    \begin{itemize}
        \item Rappelons avec la plus grande rigueur la définition de la racine carrée : pour tout réel $u \in \R$, $\sqrt{u^2} = |u|$.
        \item Ici, $e^{2x} = (e^x)^2$. Par conséquent :
        \[
        \sqrt{e^{2x}} = \sqrt{(e^x)^2} = |e^x|.
        \]
        \item Or, d'après les propriétés fondamentales du cours, la fonction exponentielle est \textbf{strictement positive} sur $\R$ : pour tout $x \in \R$, $e^x > 0$. La valeur absolue d'un nombre strictement positif est le nombre lui-même : $|e^x| = e^x$.
        \item \textit{Méthode alternative (Exposant rationnel) :}
        Pour tout $Y > 0$, $\sqrt{Y} = Y^{1/2}$. Ainsi :
        \[
        \sqrt{e^{2x}} = \left(e^{2x}\right)^{1/2} = e^{2x \times \frac{1}{2}} = e^x.
        \]
    \end{itemize}
    \textbf{Conclusion :} $\boxed{A_4(x) = e^x}$.

    \item \textbf{Simplification de $A_5(x) = \sqrt{e^{-4x}}$ :}
    \begin{itemize}
        \item Écrivons $e^{-4x}$ sous la forme d'un carré parfait : $e^{-4x} = (e^{-2x})^2$.
        \item Ainsi :
        \[
        \sqrt{e^{-4x}} = \sqrt{(e^{-2x})^2} = |e^{-2x}|.
        \]
        \item Comme pour tout réel $x$, $e^{-2x} > 0$, on a $|e^{-2x}| = e^{-2x}$.
        \item En utilisant les puissances fractionnaires :
        \[
        \sqrt{e^{-4x}} = \left(e^{-4x}\right)^{1/2} = e^{-4x \times \frac{1}{2}} = e^{-2x} = \frac{1}{e^{2x}}.
        \]
    \end{itemize}
    \textbf{Conclusion :} $\boxed{A_5(x) = e^{-2x} = \dfrac{1}{e^{2x}}}$.

    \item \textbf{Simplification de $A_6(x) = \left(\dfrac{e^x + e^{-x}}{2}\right)^2 - \left(\dfrac{e^x - e^{-x}}{2}\right)^2$ :}
    \begin{itemize}
        \item \textbf{Méthode 1 (Développement direct des carrés) :}
        Calculons séparément chaque carré :
        \[
        \left(\frac{e^x + e^{-x}}{2}\right)^2 = \frac{(e^x + e^{-x})^2}{4} = \frac{(e^x)^2 + 2 e^x e^{-x} + (e^{-x})^2}{4} = \frac{e^{2x} + 2e^0 + e^{-2x}}{4} = \frac{e^{2x} + 2 + e^{-2x}}{4}.
        \]
        De même :
        \[
        \left(\frac{e^x - e^{-x}}{2}\right)^2 = \frac{(e^x - e^{-x})^2}{4} = \frac{(e^x)^2 - 2 e^x e^{-x} + (e^{-x})^2}{4} = \frac{e^{2x} - 2e^0 + e^{-2x}}{4} = \frac{e^{2x} - 2 + e^{-2x}}{4}.
        \]
        Soustrayons les deux fractions ayant le même dénominateur 4 :
        \[
        A_6(x) = \frac{(e^{2x} + 2 + e^{-2x}) - (e^{2x} - 2 + e^{-2x})}{4} = \frac{e^{2x} + 2 + e^{-2x} - e^{2x} + 2 - e^{-2x}}{4} = \frac{4}{4} = 1.
        \]
        \item \textbf{Méthode 2 (Factorisation par identité remarquable $u^2 - v^2 = (u+v)(u-v)$) :}
        Posons $u = \dfrac{e^x + e^{-x}}{2}$ et $v = \dfrac{e^x - e^{-x}}{2}$.
        Calculons la somme et la différence :
        \[
        u + v = \frac{e^x + e^{-x} + e^x - e^{-x}}{2} = \frac{2e^x}{2} = e^x,
        \]
        \[
        u - v = \frac{e^x + e^{-x} - (e^x - e^{-x})}{2} = \frac{e^x + e^{-x} - e^x + e^{-x}}{2} = \frac{2e^{-x}}{2} = e^{-x}.
        \]
        Par le produit remarquable :
        \[
        A_6(x) = u^2 - v^2 = (u+v)(u-v) = e^x \times e^{-x} = e^{x - x} = e^0 = 1.
        \]
    \end{itemize}
    \textbf{Conclusion :} $\boxed{A_6(x) = 1}$.
\end{enumerate}
\end{correctionbox}

\begin{remarquebox}[Lien structurel avec la trigonométrie hyperbolique]
L'expression $A_6(x)$ n'est autre que la relation fondamentale de la trigonométrie hyperbolique. En posant :
\[
\ch x = \frac{e^x + e^{-x}}{2} \quad \text{(cosinus hyperbolique)} \qquad \text{et} \qquad \sh x = \frac{e^x - e^{-x}}{2} \quad \text{(sinus hyperbolique)},
\]
le résultat obtenu s'écrit de façon canonique :
\[
\ch^2 x - \sh^2 x = 1.
\]
Tandis que les fonctions trigonométriques circulaires usuelles paramètrent le cercle euclidien $X^2 + Y^2 = 1$ via $(X, Y) = (\cos t, \sin t)$, les fonctions hyperboliques paramètrent la branche droite de l'hyperbole équilatère $X^2 - Y^2 = 1$ via $(X, Y) = (\ch t, \sh t)$.
\end{remarquebox}

% =========================================================================
% EXERCICE 2
% =========================================================================
\subsection*{Exercice 2 : Preuve d'identités universelles pour tout $x \in \R$}
\addcontentsline{toc}{subsection}{Exercice 2 : Preuve d'identités universelles pour tout $x \in \R$}

Démontrer que pour tout réel $x \in \R$, les égalités suivantes sont vérifiées :
\begin{enumerate}[label=\textbf{\arabic*.}]
    \item $\dfrac{e^x - 1}{e^x + 1} = \dfrac{1 - e^{-x}}{1 + e^{-x}}$
    \item $\dfrac{e^{2x} - 1}{e^x + 1} = e^x \dfrac{1 - e^{-2x}}{1 + e^{-x}} = e^x - 1$
    \item $e^{-x} - e^{-2x} = \dfrac{e^x - 1}{e^{2x}}$
\end{enumerate}

\begin{correctionbox}[Exercice 2]
\begin{enumerate}[label=\textbf{\arabic*.}]
    \item \textbf{Démonstration de l'identité $\dfrac{e^x - 1}{e^x + 1} = \dfrac{1 - e^{-x}}{1 + e^{-x}}$ :}
    \begin{itemize}
        \item \textbf{Justification du domaine de validité :}
        Pour tout $x \in \R$, $e^x > 0$, donc $e^x + 1 > 1 > 0$. De même, $e^{-x} > 0$, donc $1 + e^{-x} > 1 > 0$. Les deux dénominateurs sont strictement positifs et ne s'annulent jamais sur $\R$. Les deux membres sont donc définis pour tout $x \in \R$.
        \item \textbf{Méthode 1 (Factorisation par $e^x$ au numérateur et au dénominateur) :}
        Partons du membre de gauche :
        \[
        \frac{e^x - 1}{e^x + 1} = \frac{e^x\left(1 - \frac{1}{e^x}\right)}{e^x\left(1 + \frac{1}{e^x}\right)}.
        \]
        Puisque $e^x \neq 0$, nous pouvons simplifier par $e^x$ :
        \[
        \frac{e^x - 1}{e^x + 1} = \frac{1 - \frac{1}{e^x}}{1 + \frac{1}{e^x}}.
        \]
        Comme $\dfrac{1}{e^x} = e^{-x}$, on obtient immédiatement :
        \[
        \frac{e^x - 1}{e^x + 1} = \frac{1 - e^{-x}}{1 + e^{-x}}.
        \]
        \item \textbf{Méthode 2 (Multiplication par $e^{-x}$) :}
        Multiplions numérateur et dénominateur du membre de gauche par $e^{-x} \neq 0$ :
        \[
        \frac{(e^x - 1)e^{-x}}{(e^x + 1)e^{-x}} = \frac{e^x e^{-x} - e^{-x}}{e^x e^{-x} + e^{-x}} = \frac{e^0 - e^{-x}}{e^0 + e^{-x}} = \frac{1 - e^{-x}}{1 + e^{-x}}.
        \]
        L'égalité est démontrée.
    \end{itemize}

    \item \textbf{Démonstration de la chaîne d'égalités $\dfrac{e^{2x} - 1}{e^x + 1} = e^x \dfrac{1 - e^{-2x}}{1 + e^{-x}} = e^x - 1$ :}
    \begin{itemize}
        \item \textbf{Étape 1 : Simplification de $\dfrac{e^{2x} - 1}{e^x + 1}$ :}
        Remarquons que le numérateur est une différence de deux carrés :
        \[
        e^{2x} - 1 = (e^x)^2 - 1^2 = (e^x - 1)(e^x + 1).
        \]
        Le dénominateur $e^x + 1$ étant strictement positif pour tout $x \in \R$ (donc non nul), on peut simplifier la fraction :
        \[
        \frac{e^{2x} - 1}{e^x + 1} = \frac{(e^x - 1)(e^x + 1)}{e^x + 1} = e^x - 1.
        \]
        \item \textbf{Étape 2 : Simplification du terme médian $e^x \dfrac{1 - e^{-2x}}{1 + e^{-x}}$ :}
        Développons le numérateur en distribuant $e^x$ :
        \[
        e^x (1 - e^{-2x}) = e^x - e^x e^{-2x} = e^x - e^{x - 2x} = e^x - e^{-x}.
        \]
        Or, par factorisation d'une différence de carrés en posant $A = 1$ et $B = e^{-x}$ :
        \[
        1 - e^{-2x} = 1^2 - (e^{-x})^2 = (1 - e^{-x})(1 + e^{-x}).
        \]
        Par conséquent, dans l'expression médiane, le facteur $(1 + e^{-x}) > 0$ se simplifie :
        \[
        e^x \frac{1 - e^{-2x}}{1 + e^{-x}} = e^x \frac{(1 - e^{-x})(1 + e^{-x})}{1 + e^{-x}} = e^x (1 - e^{-x}) = e^x - e^x e^{-x} = e^x - e^0 = e^x - 1.
        \]
        \item Les deux expressions étant rigoureusement égales à $e^x - 1$, la double égalité est établie pour tout $x \in \R$.
    \end{itemize}

    \item \textbf{Démonstration de l'identité $e^{-x} - e^{-2x} = \dfrac{e^x - 1}{e^{2x}}$ :}
    \begin{itemize}
        \item \textbf{Méthode 1 (En partant du membre de droite) :}
        Le dénominateur $e^{2x}$ est strictement positif. Décomposons la fraction en différence de deux termes :
        \[
        \frac{e^x - 1}{e^{2x}} = \frac{e^x}{e^{2x}} - \frac{1}{e^{2x}}.
        \]
        En appliquant les propriétés du quotient :
        \[
        \frac{e^x}{e^{2x}} = e^{x - 2x} = e^{-x} \qquad \text{et} \qquad \frac{1}{e^{2x}} = e^{-2x}.
        \]
        On retrouve directement :
        \[
        \frac{e^x - 1}{e^{2x}} = e^{-x} - e^{-2x}.
        \]
        \item \textbf{Méthode 2 (En partant du membre de gauche par factorisation) :}
        Dans l'expression $e^{-x} - e^{-2x}$, factorisons par le terme de plus grand exposant négatif $e^{-2x}$ :
        \[
        e^{-x} - e^{-2x} = e^{-2x}\left(\frac{e^{-x}}{e^{-2x}} - 1\right) = e^{-2x}\left(e^{-x - (-2x)} - 1\right) = e^{-2x}(e^x - 1).
        \]
        Puisque $e^{-2x} = \dfrac{1}{e^{2x}}$, il vient immédiatement :
        \[
        e^{-x} - e^{-2x} = \frac{e^x - 1}{e^{2x}}.
        \]
    \end{itemize}
\end{enumerate}
\end{correctionbox}

\begin{methodebox}[Technique de l'arc-moitié / Exponentielle médiane (Réflexe CPGE)]
Dans de nombreux calculs trigonométriques ou exponentiels avancés, la factorisation par l'\textbf{exponentielle à mi-pente} (ou angle moitié) simplifie considérablement les expressions et fait apparaître les fonctions hyperboliques :
\[
e^x - 1 = e^{x/2}\left(e^{x/2} - e^{-x/2}\right) = 2e^{x/2} \sh(x/2),
\]
\[
e^x + 1 = e^{x/2}\left(e^{x/2} + e^{-x/2}\right) = 2e^{x/2} \ch(x/2).
\]
En reprenant l'identité 1 de l'Exercice 2, on a :
\[
\frac{e^x - 1}{e^x + 1} = \frac{2e^{x/2} \sh(x/2)}{2e^{x/2} \ch(x/2)} = \frac{\sh(x/2)}{\ch(x/2)} = \th\left(\frac{x}{2}\right).
\]
Cette factorisation systématique par l'argument médian est l'un des réflexes les plus puissants en Mathématiques Supérieures pour lever des indéterminations ou simplifier des sommes trigonométriques.
\end{methodebox}

\vspace{0.8cm}
\hrule
\vspace{0.8cm}
\section{Section 2 : Dérivation et Études de variations de base}

L'association de la dérivation et de la fonction exponentielle ouvre la voie à l'étude fine des variations et du comportement asymptotique de classes fonctionnelles très variées (fonctions gaussiennes, fonctions puissances pondérées, quotients d'exponentielles). L'enjeu central réside dans la factorisation systématique de la dérivée par l'exponentielle, dont la stricte positivité permet de ramener l'étude du signe de $f'(x)$ à celle d'un facteur polynomial ou rationnel usuel.

% =========================================================================
% EXERCICE 3
% =========================================================================
\subsection*{Exercice 3 : Variations complètes des six fonctions de référence}
\addcontentsline{toc}{subsection}{Exercice 3 : Variations complètes des six fonctions de référence}

Pour chacune des six fonctions suivantes, déterminer son ensemble de définition, justifier sa dérivabilité, expliciter sa dérivée, étudier rigoureusement son signe, et dresser son tableau de variations complet (incluant les extrema éventuels et les limites aux bornes) :
\begin{enumerate}[label=\textbf{\arabic*.}]
    \item $f_1(x) = e^{-x}$
    \item $f_2(x) = x - e^x$
    \item $f_3(x) = e^x + e^{-x}$
    \item $f_4(x) = e^{3x} - 3e^x$
    \item $f_5(x) = e^{-x^2}$
    \item $f_6(x) = (x+1)^2 e^{-x}$
\end{enumerate}

\begin{correctionbox}[Exercice 3]
Traitons exhaustivement chaque fonction selon le protocole d'analyse canonique.

% ----------------------------------------------------
% f_1
% ----------------------------------------------------
\subsubsection*{1. Étude de $f_1(x) = e^{-x}$}
\begin{itemize}
    \item \textbf{Ensemble de définition et dérivabilité :}
    La fonction affine $u_1 : x \mapsto -x$ est dérivable sur $\R$ et $\exp$ est dérivable sur $\R$. Par composition, $f_1 = \exp \circ u_1$ est dérivable sur $\mathcal{D}_{f_1} = \R$.
    \item \textbf{Calcul de la dérivée :}
    En appliquant la formule $(e^u)' = u' e^u$ avec $u_1'(x) = -1$ :
    \[
    \forall x \in \R, \quad f_1'(x) = -e^{-x}.
    \]
    \item \textbf{Étude du signe et sens de variation :}
    Pour tout réel $x$, $e^{-x} > 0$, donc $f_1'(x) = -e^{-x} < 0$. La dérivée est strictement négative sur $\R$. Par suite, $f_1$ est \textbf{strictement décroissante} sur $\R$.
    \item \textbf{Limites aux bornes :}
    \begin{itemize}
        \item En $-\infty$ : Lorsque $x \to -\infty$, $u_1(x) = -x \to +\infty$. Par composition avec $\lim_{X \to +\infty} e^X = +\infty$, on a $\lim_{x \to -\infty} f_1(x) = +\infty$.
        \item En $+\infty$ : Lorsque $x \to +\infty$, $u_1(x) = -x \to -\infty$. Par composition avec $\lim_{X \to -\infty} e^X = 0$, on a $\lim_{x \to +\infty} f_1(x) = 0$.
    \end{itemize}
    L'axe des abscisses ($y = 0$) est asymptote horizontale à la courbe en $+\infty$.
    \item \textbf{Tableau de variations :}
    \begin{center}
    \begin{tabular}{|c|lcccr|}
    \hline
    $x$ & $-\infty$ & & & & $+\infty$ \\
    \hline
    $f_1'(x)$ & & & $-$ & & \\
    \hline
    $f_1(x)$ & $+\infty$ & & $\searrow$ & & $0$ \\
    \hline
    \end{tabular}
    \end{center}
\end{itemize}

% ----------------------------------------------------
% f_2
% ----------------------------------------------------
\subsubsection*{2. Étude de $f_2(x) = x - e^x$}
\begin{itemize}
    \item \textbf{Ensemble de définition et dérivabilité :}
    $f_2$ est la somme de la fonction identité $x \mapsto x$ et de $x \mapsto -e^x$, toutes deux dérivables sur $\R$. Donc $f_2$ est dérivable sur $\mathcal{D}_{f_2} = \R$.
    \item \textbf{Calcul de la dérivée :}
    \[
    \forall x \in \R, \quad f_2'(x) = 1 - e^x.
    \]
    \item \textbf{Étude du signe de $f_2'(x)$ :}
    Résolvons l'équation et l'inéquation associées :
    \[
    f_2'(x) = 0 \iff 1 - e^x = 0 \iff e^x = 1 \iff x = 0,
    \]
    \[
    f_2'(x) > 0 \iff 1 - e^x > 0 \iff e^x < 1 \iff x < 0 \quad \text{(par stricte croissance de $\exp$)},
    \]
    \[
    f_2'(x) < 0 \iff 1 - e^x < 0 \iff e^x > 1 \iff x > 0.
    \]
    \item \textbf{Variations et extremum :}
    $f_2$ est strictement croissante sur $]-\infty, 0]$ et strictement décroissante sur $[0, +\infty[$.
    Elle admet en $x = 0$ un \textbf{maximum absolu strict} qui vaut :
    \[
    f_2(0) = 0 - e^0 = -1.
    \]
    \item \textbf{Conséquence fondamentale :}
    Puisque $-1$ est le maximum absolu de $f_2$ sur $\R$, on a pour tout $x \in \R$ :
    \[
    f_2(x) \le -1 < 0 \iff x - e^x < 0 \iff e^x > x.
    \]
    Cette inégalité fondamentale garantit que la courbe de l'exponentielle est strictement au-dessus de la première bissectrice sur tout $\R$.
    \item \textbf{Limites aux bornes :}
    \begin{itemize}
        \item En $-\infty$ : $\lim_{x \to -\infty} x = -\infty$ et $\lim_{x \to -\infty} e^x = 0$, donc par somme : $\lim_{x \to -\infty} f_2(x) = -\infty$. De plus, $f_2(x) - x = -e^x \to 0$, donc la droite $y = x$ est asymptote oblique à la courbe en $-\infty$.
        \item En $+\infty$ : C'est une forme indéterminée « $\infty - \infty$ ». Factorisons par le terme prépondérant $e^x$ :
        \[
        f_2(x) = e^x \left(\frac{x}{e^x} - 1\right).
        \]
        Par croissances comparées, $\lim_{x \to +\infty} \frac{x}{e^x} = 0$. Donc $\lim_{x \to +\infty} \left(\frac{x}{e^x} - 1\right) = -1$.
        Comme $\lim_{x \to +\infty} e^x = +\infty$, par produit : $\lim_{x \to +\infty} f_2(x) = -\infty$.
    \end{itemize}
    \item \textbf{Tableau de variations :}
    \begin{center}
    \begin{tabular}{|c|lccccr|}
    \hline
    $x$ & $-\infty$ & & $0$ & & $+\infty$ \\
    \hline
    $f_2'(x)$ & & $+$ & $0$ & $-$ & \\
    \hline
    $f_2(x)$ & $-\infty$ & $\nearrow$ & $-1$ & $\searrow$ & $-\infty$ \\
    \hline
    \end{tabular}
    \end{center}
\end{itemize}

% ----------------------------------------------------
% f_3
% ----------------------------------------------------
\subsubsection*{3. Étude de $f_3(x) = e^x + e^{-x}$}
\begin{itemize}
    \item \textbf{Propriétés de symétrie (Parité) :}
    Pour tout $x \in \R$, $-x \in \R$ et :
    \[
    f_3(-x) = e^{-x} + e^{-(-x)} = e^{-x} + e^x = f_3(x).
    \]
    La fonction $f_3$ est \textbf{paire}. Sa courbe représentative admet l'axe des ordonnées $(Oy)$ pour axe de symétrie orthogonal. On peut restreindre son étude à l'intervalle $[0, +\infty[$.
    \item \textbf{Dérivabilité et calcul de dérivée :}
    $f_3$ est dérivable sur $\R$ comme somme de fonctions dérivables :
    \[
    \forall x \in \R, \quad f_3'(x) = e^x - e^{-x}.
    \]
    \item \textbf{Signe de $f_3'(x)$ :}
    \[
    f_3'(x) > 0 \iff e^x - e^{-x} > 0 \iff e^x > e^{-x} \iff x > -x \iff 2x > 0 \iff x > 0.
    \]
    De même, $f_3'(x) = 0 \iff x = 0$, et $f_3'(x) < 0 \iff x < 0$.
    \item \textbf{Variations et extremum :}
    $f_3$ est strictement décroissante sur $]-\infty, 0]$ et strictement croissante sur $[0, +\infty[$.
    Elle admet en $x = 0$ un \textbf{minimum absolu strict} valant :
    \[
    f_3(0) = e^0 + e^0 = 1 + 1 = 2.
    \]
    Pour tout $x \in \R$, $e^x + e^{-x} \ge 2$, avec égalité si et seulement si $x = 0$.
    \item \textbf{Limites :}
    Par parité, $\lim_{x \to -\infty} f_3(x) = \lim_{x \to +\infty} f_3(x)$.
    Comme $\lim_{x \to +\infty} e^x = +\infty$ et $\lim_{x \to +\infty} e^{-x} = 0$, on a $\lim_{x \to +\infty} f_3(x) = +\infty$ (et de même en $-\infty$).
    \item \textbf{Tableau de variations :}
    \begin{center}
    \begin{tabular}{|c|lccccr|}
    \hline
    $x$ & $-\infty$ & & $0$ & & $+\infty$ \\
    \hline
    $f_3'(x)$ & & $-$ & $0$ & $+$ & \\
    \hline
    $f_3(x)$ & $+\infty$ & $\searrow$ & $2$ & $\nearrow$ & $+\infty$ \\
    \hline
    \end{tabular}
    \end{center}
\end{itemize}

% ----------------------------------------------------
% f_4
% ----------------------------------------------------
\subsubsection*{4. Étude de $f_4(x) = e^{3x} - 3e^x$}
\begin{itemize}
    \item \textbf{Dérivabilité et calcul de la dérivée :}
    $f_4$ est dérivable sur $\R$ comme combinaison linéaire de fonctions exponentielles composées :
    \[
    f_4'(x) = 3e^{3x} - 3e^x = 3e^x (e^{2x} - 1).
    \]
    \item \textbf{Étude du signe de $f_4'(x)$ :}
    Pour tout $x \in \R$, $3e^x > 0$. Le signe de $f_4'(x)$ est donc rigoureusement identique à celui du facteur $(e^{2x} - 1)$ :
    \[
    e^{2x} - 1 = 0 \iff e^{2x} = 1 \iff 2x = 0 \iff x = 0,
    \]
    \[
    e^{2x} - 1 > 0 \iff e^{2x} > 1 \iff 2x > 0 \iff x > 0,
    \]
    \[
    e^{2x} - 1 < 0 \iff e^{2x} < 1 \iff 2x < 0 \iff x < 0.
    \]
    \item \textbf{Variations et extremum :}
    $f_4$ est strictement décroissante sur $]-\infty, 0]$ et strictement croissante sur $[0, +\infty[$.
    Elle atteint son \textbf{minimum absolu strict} en $x = 0$ :
    \[
    f_4(0) = e^0 - 3e^0 = 1 - 3 = -2.
    \]
    \item \textbf{Limites aux bornes :}
    \begin{itemize}
        \item En $-\infty$ : $\lim_{x \to -\infty} 3x = -\infty \implies \lim_{x \to -\infty} e^{3x} = 0$, et $\lim_{x \to -\infty} 3e^x = 0$. Par somme, $\lim_{x \to -\infty} f_4(x) = 0$.
        L'axe des abscisses est asymptote horizontale en $-\infty$.
        \item En $+\infty$ : Forme indéterminée « $\infty - \infty$ ». Factorisons par le terme prépondérant $e^{3x}$ :
        \[
        f_4(x) = e^{3x}\left(1 - 3\frac{e^x}{e^{3x}}\right) = e^{3x}(1 - 3e^{-2x}).
        \]
        Comme $\lim_{x \to +\infty} e^{-2x} = 0$, la parenthèse tend vers $1$. Comme $\lim_{x \to +\infty} e^{3x} = +\infty$, on conclut par produit : $\lim_{x \to +\infty} f_4(x) = +\infty$.
    \end{itemize}
    \item \textbf{Tableau de variations :}
    \begin{center}
    \begin{tabular}{|c|lccccr|}
    \hline
    $x$ & $-\infty$ & & $0$ & & $+\infty$ \\
    \hline
    $f_4'(x)$ & & $-$ & $0$ & $+$ & \\
    \hline
    $f_4(x)$ & $0$ & $\searrow$ & $-2$ & $\nearrow$ & $+\infty$ \\
    \hline
    \end{tabular}
    \end{center}
\end{itemize}

% ----------------------------------------------------
% f_5
% ----------------------------------------------------
\subsubsection*{5. Étude de $f_5(x) = e^{-x^2}$}
\begin{itemize}
    \item \textbf{Parité et régularité :}
    Pour tout $x \in \R$, $f_5(-x) = e^{-(-x)^2} = e^{-x^2} = f_5(x)$. La fonction est \textbf{paire}.
    De plus, le polynôme $u_5(x) = -x^2$ est dérivable sur $\R$, donc $f_5$ est dérivable sur $\R$.
    \item \textbf{Calcul de la dérivée :}
    En appliquant la formule $(e^u)' = u' e^u$ avec $u_5'(x) = -2x$ :
    \[
    \forall x \in \R, \quad f_5'(x) = -2x e^{-x^2}.
    \]
    \item \textbf{Étude du signe de $f_5'(x)$ :}
    Puisque pour tout $x \in \R$, $e^{-x^2} > 0$, le signe de $f_5'(x)$ est exactement le signe opposé de $x$ :
    \[
    f_5'(x) > 0 \iff x < 0, \qquad f_5'(0) = 0, \qquad f_5'(x) < 0 \iff x > 0.
    \]
    \item \textbf{Variations et extremum :}
    $f_5$ est strictement croissante sur $]-\infty, 0]$ et strictement décroissante sur $[0, +\infty[$.
    Elle admet en $x = 0$ un \textbf{maximum absolu strict} valant :
    \[
    f_5(0) = e^0 = 1.
    \]
    \item \textbf{Limites aux bornes :}
    Comme $\lim_{x \to \pm\infty} (-x^2) = -\infty$ et $\lim_{X \to -\infty} e^X = 0$, on a par composition :
    \[
    \lim_{x \to -\infty} f_5(x) = 0 \qquad \text{et} \qquad \lim_{x \to +\infty} f_5(x) = 0.
    \]
    L'axe des abscisses ($y = 0$) est asymptote horizontale à la courbe en $-\infty$ et en $+\infty$.
    \item \textbf{Tableau de variations :}
    \begin{center}
    \begin{tabular}{|c|lccccr|}
    \hline
    $x$ & $-\infty$ & & $0$ & & $+\infty$ \\
    \hline
    $f_5'(x)$ & & $+$ & $0$ & $-$ & \\
    \hline
    $f_5(x)$ & $0$ & $\nearrow$ & $1$ & $\searrow$ & $0$ \\
    \hline
    \end{tabular}
    \end{center}
\end{itemize}

% ----------------------------------------------------
% f_6
% ----------------------------------------------------
\subsubsection*{6. Étude de $f_6(x) = (x+1)^2 e^{-x}$}
\begin{itemize}
    \item \textbf{Dérivabilité et calcul de dérivée :}
    $f_6$ est le produit des fonctions $u(x) = (x+1)^2$ et $v(x) = e^{-x}$, toutes deux dérivables sur $\R$.
    Leurs dérivées respectives sont $u'(x) = 2(x+1)$ et $v'(x) = -e^{-x}$. Par la règle du produit $(uv)' = u'v + uv'$ :
    \[
    f_6'(x) = 2(x+1)e^{-x} + (x+1)^2 (-e^{-x}) = (x+1)\bigl[2 - (x+1)\bigr]e^{-x} = (x+1)(1 - x)e^{-x}.
    \]
    En développant le produit des facteurs affines :
    \[
    (x+1)(1-x) = 1 - x^2.
    \]
    D'où la forme compacte de la dérivée :
    \[
    \forall x \in \R, \quad f_6'(x) = (1 - x^2)e^{-x}.
    \]
    \item \textbf{Étude du signe de $f_6'(x)$ :}
    Pour tout $x \in \R$, $e^{-x} > 0$. Le signe de $f_6'(x)$ est donc rigoureusement identique à celui du trinôme du second degré $1 - x^2$.
    Les racines de $1 - x^2 = 0$ sont $x = -1$ et $x = 1$.
    Le coefficient devant $x^2$ étant négatif ($-1 < 0$), la parabole est orientée vers le bas :
    \begin{itemize}
        \item $1 - x^2 < 0 \iff x \in ]-\infty, -1[ \cup ]1, +\infty[$,
        \item $1 - x^2 = 0 \iff x \in \{-1, 1\}$,
        \item $1 - x^2 > 0 \iff x \in ]-1, 1[$.
    \end{itemize}
    \item \textbf{Variations et extrema :}
    \begin{itemize}
        \item $f_6$ est strictement décroissante sur $]-\infty, -1]$,
        \item strictement croissante sur $[-1, 1]$,
        \item strictement décroissante sur $[1, +\infty[$.
    \end{itemize}
    Extrema :
    \begin{itemize}
        \item En $x = -1$, la fonction admet un \textbf{minimum local et global} qui vaut $f_6(-1) = (-1+1)^2 e^1 = 0$ (valeur minimale car un carré est toujours positif ou nul).
        \item En $x = 1$, la fonction admet un \textbf{maximum local} qui vaut $f_6(1) = (1+1)^2 e^{-1} = 4e^{-1} = \dfrac{4}{e} \approx 1{,}4715$.
    \end{itemize}
    \item \textbf{Limites aux bornes :}
    \begin{itemize}
        \item En $-\infty$ : $\lim_{x \to -\infty} (x+1)^2 = +\infty$ et $\lim_{x \to -\infty} e^{-x} = +\infty$. Par produit : $\lim_{x \to -\infty} f_6(x) = +\infty$.
        \item En $+\infty$ : Forme indéterminée « $\infty \times 0$ ». Développons l'expression :
        \[
        f_6(x) = (x^2 + 2x + 1)e^{-x} = x^2 e^{-x} + 2x e^{-x} + e^{-x}.
        \]
        D'après les théorèmes de croissances comparées en $+\infty$ :
        \[
        \lim_{x \to +\infty} x^2 e^{-x} = 0, \qquad \lim_{x \to +\infty} x e^{-x} = 0, \qquad \lim_{x \to +\infty} e^{-x} = 0.
        \]
        Par somme de limites finies, on obtient :
        \[
        \lim_{x \to +\infty} f_6(x) = 0.
        \]
        La courbe $\mathcal{C}_{f_6}$ admet donc l'axe des abscisses pour asymptote horizontale en $+\infty$.
    \end{itemize}
    \item \textbf{Tableau de variations :}
    \begin{center}
    \begin{tabular}{|c|lcccccccr|}
    \hline
    $x$ & $-\infty$ & & $-1$ & & $1$ & & $+\infty$ \\
    \hline
    $f_6'(x)$ & & $-$ & $0$ & $+$ & $0$ & $-$ & \\
    \hline
    $f_6(x)$ & $+\infty$ & $\searrow$ & $0$ & $\nearrow$ & $\dfrac{4}{e}$ & $\searrow$ & $0$ \\
    \hline
    \end{tabular}
    \end{center}
\end{itemize}
\end{correctionbox}

% =========================================================================
% EXERCICE 4
% =========================================================================
\subsection*{Exercice 4 : Calculs de dérivées détaillées et formes optimales}
\addcontentsline{toc}{subsection}{Exercice 4 : Calculs de dérivées détaillées et formes optimales}

Pour chacune des fonctions suivantes, préciser son domaine de définition $\mathcal{D}_f$, son ensemble de dérivabilité, et déterminer l'expression de sa dérivée sous la forme la plus simplifiée et factorisée possible :
\begin{enumerate}[label=\textbf{\arabic*.}]
    \item $f_1(x) = \dfrac{2e^x - 1}{e^x + 3}$
    \item $f_2(x) = e^{\frac{2x+1}{x+4}}$
    \item $f_3(x) = x e^{-x^2}$
    \item $f_4(x) = e^{1/x}$
    \item $f_5(x) = \dfrac{3}{1 + e^{2x}}$
    \item $f_6(x) = x + 1 - 3x e^{-x^2}$
\end{enumerate}

\begin{correctionbox}[Exercice 4]
\begin{enumerate}[label=\textbf{\arabic*.}]
    \item \textbf{Dérivée de $f_1(x) = \dfrac{2e^x - 1}{e^x + 3}$ :}
    \begin{itemize}
        \item \textbf{Domaine :} Pour tout $x \in \R$, $e^x > 0 \implies e^x + 3 > 3 > 0$. Le dénominateur ne s'annule jamais. $f_1$ est définie et dérivable sur $\R$ comme quotient de fonctions dérivables.
        \item \textbf{Formule du quotient :} Posons $u(x) = 2e^x - 1 \implies u'(x) = 2e^x$ et $v(x) = e^x + 3 \implies v'(x) = e^x$.
        \[
        f_1'(x) = \frac{u'(x)v(x) - u(x)v'(x)}{\bigl(v(x)\bigr)^2} = \frac{2e^x(e^x + 3) - (2e^x - 1)e^x}{(e^x + 3)^2}.
        \]
        \item \textbf{Simplification du numérateur :}
        \[
        2e^x(e^x + 3) - (2e^x - 1)e^x = 2e^{2x} + 6e^x - 2e^{2x} + e^x = 7e^x.
        \]
        \item \textbf{Conclusion :}
        \[
        \boxed{f_1'(x) = \frac{7e^x}{(e^x + 3)^2}}.
        \]
        \textit{Remarque :} Comme $7e^x > 0$ et $(e^x+3)^2 > 0$, $f_1'(x) > 0$ sur $\R$ : $f_1$ est strictement croissante sur $\R$.
    \end{itemize}

    \item \textbf{Dérivée de $f_2(x) = e^{\frac{2x+1}{x+4}}$ :}
    \begin{itemize}
        \item \textbf{Domaine :} La fonction homographique $u(x) = \dfrac{2x+1}{x+4}$ est définie et dérivable sur $\R \setminus \{-4\}$. Par composition avec l'exponentielle, $f_2$ est définie et dérivable sur $\mathcal{D}_{f_2} = ]-\infty, -4[ \cup ]-4, +\infty[$.
        \item \textbf{Dérivée de l'exposant $u(x)$ :}
        \[
        u'(x) = \frac{2(x+4) - (2x+1)(1)}{(x+4)^2} = \frac{2x + 8 - 2x - 1}{(x+4)^2} = \frac{7}{(x+4)^2}.
        \]
        \item \textbf{Dérivée de la composée :} Par la formule $\bigl(e^u\bigr)' = u' e^u$ :
        \[
        \boxed{f_2'(x) = \frac{7}{(x+4)^2} \exp\left(\frac{2x+1}{x+4}\right)}.
        \]
        \textit{Remarque :} Pour tout $x \neq -4$, $f_2'(x) > 0$. La fonction est donc strictement croissante sur chacun des deux intervalles $]-\infty, -4[$ et $]-4, +\infty[$.
    \end{itemize}

    \item \textbf{Dérivée de $f_3(x) = x e^{-x^2}$ :}
    \begin{itemize}
        \item \textbf{Domaine :} Produit de fonctions dérivables sur $\R$, donc dérivable sur $\R$.
        \item \textbf{Formule du produit :} Posons $u(x) = x \implies u'(x) = 1$ et $v(x) = e^{-x^2} \implies v'(x) = -2x e^{-x^2}$.
        \[
        f_3'(x) = 1 \cdot e^{-x^2} + x \cdot \bigl(-2x e^{-x^2}\bigr) = e^{-x^2} - 2x^2 e^{-x^2}.
        \]
        \item \textbf{Factorisation optimale :}
        \[
        \boxed{f_3'(x) = (1 - 2x^2) e^{-x^2}}.
        \]
        \textit{Remarque :} Le signe de $f_3'(x)$ dépend exclusivement du trinôme $1 - 2x^2 = (1 - x\sqrt{2})(1 + x\sqrt{2})$, qui s'annule en $\pm \frac{\sqrt{2}}{2}$.
    \end{itemize}

    \item \textbf{Dérivée de $f_4(x) = e^{1/x}$ :}
    \begin{itemize}
        \item \textbf{Domaine :} La fonction inverse $x \mapsto 1/x$ est dérivable sur $\R^* = ]-\infty, 0[ \cup ]0, +\infty[$. Par composition, $f_4$ est dérivable sur $\R^*$.
        \item \textbf{Formule de la composée :} Avec $u(x) = \dfrac{1}{x} \implies u'(x) = -\dfrac{1}{x^2}$ :
        \[
        \boxed{f_4'(x) = -\frac{1}{x^2} e^{1/x}}.
        \]
        \textit{Remarque :} Pour tout $x \in \R^*$, $-\dfrac{1}{x^2} < 0$ et $e^{1/x} > 0$, donc $f_4'(x) < 0$. $f_4$ est strictement décroissante sur $]-\infty, 0[$ et strictement décroissante sur $]0, +\infty[$.
    \end{itemize}

    \item \textbf{Dérivée de $f_5(x) = \dfrac{3}{1 + e^{2x}}$ :}
    \begin{itemize}
        \item \textbf{Domaine :} Comme $e^{2x} > 0$, le dénominateur $1 + e^{2x} \ge 1 > 0$ ne s'annule jamais. $f_5$ est dérivable sur $\R$.
        \item \textbf{Formule de l'inverse :} Pour une fonction $v$ ne s'annulant pas, $\left(\dfrac{1}{v}\right)' = -\dfrac{v'}{v^2}$.
        Ici $v(x) = 1 + e^{2x} \implies v'(x) = 2e^{2x}$. Par linéarité :
        \[
        f_5'(x) = 3 \times \left( -\frac{2e^{2x}}{(1 + e^{2x})^2} \right) = \frac{-6e^{2x}}{(1 + e^{2x})^2}.
        \]
        \item \textbf{Conclusion :}
        \[
        \boxed{f_5'(x) = -\frac{6e^{2x}}{(1 + e^{2x})^2}}.
        \]
        \textit{Remarque :} La dérivée est strictement négative sur $\R$, donc $f_5$ est strictement décroissante sur $\R$.
    \end{itemize}

    \item \textbf{Dérivée de $f_6(x) = x + 1 - 3x e^{-x^2}$ :}
    \begin{itemize}
        \item \textbf{Domaine :} Somme et produit de fonctions usuelles dérivables sur $\R$, donc dérivable sur $\R$.
        \item \textbf{Décomposition du calcul :}
        La dérivée de $x \mapsto x + 1$ est $1$.
        D'après la question 3, la dérivée de $x \mapsto x e^{-x^2}$ est $(1 - 2x^2)e^{-x^2}$.
        Par linéarité de la dérivation :
        \[
        f_6'(x) = 1 - 3 \times \Bigl[(1 - 2x^2)e^{-x^2}\Bigr] = 1 - 3(1 - 2x^2)e^{-x^2}.
        \]
        En réordonnant les termes :
        \[
        1 - 3(1 - 2x^2)e^{-x^2} = 1 + 3(2x^2 - 1)e^{-x^2} = 1 + (6x^2 - 3)e^{-x^2}.
        \]
        \item \textbf{Conclusion :}
        \[
        \boxed{f_6'(x) = 1 + 3(2x^2 - 1)e^{-x^2} = 1 + (6x^2 - 3)e^{-x^2}}.
        \]
    \end{itemize}
\end{enumerate}
\end{correctionbox}

\vspace{0.8cm}
\hrule
\vspace{0.8cm}
\section{Section 3 : Études de fonctions, Tangentes et Familles de courbes}

Cette section rassemble les problèmes les plus stimulants et géométriquement profonds de la feuille de travaux dirigés. L'interaction entre calcul infinitésimal (tangentes, conditions de contact, enveloppes) et analyse fonctionnelle (familles de courbes à paramètres, points fixes universels, bifurcations de sens de variation) y est explorée dans l'esprit direct des épreuves de concours des grandes écoles.

% =========================================================================
% EXERCICE 5
% =========================================================================
\subsection*{Exercice 5 : Tangente passant par l'origine à une courbe exponentielle}
\addcontentsline{toc}{subsection}{Exercice 5 : Tangente passant par l'origine}

Soit $f$ la fonction définie sur $\R$ par :
\[
f(x) = x e^{x-1} + 1.
\]
On note $\mathcal{C}_f$ sa courbe représentative dans un repère orthonormé $(O; \vec{i}, \vec{j})$.
\begin{enumerate}[label=\textbf{\arabic*.}]
    \item Justifier la dérivabilité de $f$ sur $\R$ et exprimer $f'(x)$ pour tout réel $x$.
    \item Pour tout réel $a \in \R$, déterminer l'équation réduite de la tangente $T_a$ à $\mathcal{C}_f$ au point d'abscisse $a$.
    \item Montrer que la tangente $T_a$ passe par l'origine du repère $O(0, 0)$ si et seulement si :
    \[
    a^2 e^{a-1} = 1.
    \]
    \item Résoudre rigoureusement cette équation sur $\R$ en démontrant qu'elle admet $a = 1$ pour unique solution.
    \item En déduire l'équation de cette tangente remarquable et illustrer la configuration géométrique par une figure soignée.
\end{enumerate}

\begin{correctionbox}[Exercice 5]
\begin{enumerate}[label=\textbf{\arabic*.}]
    \item \textbf{Dérivabilité et calcul de $f'(x)$ :}
    La fonction $x \mapsto x-1$ est dérivable sur $\R$, donc par composition avec l'exponentielle, $x \mapsto e^{x-1}$ est dérivable sur $\R$.
    Par produit et somme avec la constante 1, $f$ est dérivable sur $\R$.
    Appliquons la règle du produit $(uv)' = u'v + uv'$ à $u(x) = x \implies u'(x) = 1$ et $v(x) = e^{x-1} \implies v'(x) = 1 \cdot e^{x-1} = e^{x-1}$ :
    \[
    f'(x) = 1 \cdot e^{x-1} + x \cdot e^{x-1} + 0 = (x+1)e^{x-1}.
    \]
    Pour tout réel $x \in \R$, $\boxed{f'(x) = (x+1)e^{x-1}}$.

    \item \textbf{Équation réduite de la tangente $T_a$ :}
    La tangente $T_a$ au point de coordonnées $(a, f(a))$ a pour équation :
    \[
    T_a : y = f'(a)(x - a) + f(a).
    \]
    Exprimons $f(a)$ et $f'(a)$ :
    \[
    f(a) = a e^{a-1} + 1 \qquad \text{et} \qquad f'(a) = (a+1)e^{a-1}.
    \]
    Développons l'équation de la droite :
    \[
    y = (a+1)e^{a-1} x - a(a+1)e^{a-1} + a e^{a-1} + 1.
    \]
    Regroupons les termes indépendants de $x$ (ordonnée à l'origine) :
    \[
    - a(a+1)e^{a-1} + a e^{a-1} + 1 = (-a^2 - a + a)e^{a-1} + 1 = -a^2 e^{a-1} + 1.
    \]
    L'équation réduite de $T_a$ s'écrit donc de façon canonique :
    \[
    \boxed{T_a : y = (a+1)e^{a-1} x + \bigl(1 - a^2 e^{a-1}\bigr)}.
    \]

    \item \textbf{Condition de passage par l'origine $O(0, 0)$ :}
    La droite $T_a$ passe par l'origine si et seulement si les coordonnées $(0, 0)$ vérifient son équation :
    \[
    0 = (a+1)e^{a-1} \times 0 + \bigl(1 - a^2 e^{a-1}\bigr) \iff 1 - a^2 e^{a-1} = 0 \iff a^2 e^{a-1} = 1.
    \]

    \item \textbf{Résolution et preuve d'unicité de la solution :}
    \begin{itemize}
        \item \textbf{Solution évidente :}
        Testons la valeur entière $a = 1$ :
        \[
        1^2 \times e^{1-1} = 1 \times e^0 = 1 \times 1 = 1.
        \]
        Donc $a = 1$ est bien une solution de l'équation.
        \item \textbf{Preuve rigoureuse de l'unicité sur $\R$ :}
        Considérons la fonction auxiliaire $\varphi : \R \to \R$ définie par :
        \[
        \varphi(x) = x^2 e^{x-1}.
        \]
        $\varphi$ est dérivable sur $\R$ comme produit de fonctions dérivables. Calculons sa dérivée :
        \[
        \varphi'(x) = 2x e^{x-1} + x^2 e^{x-1} = x(x+2)e^{x-1}.
        \]
        Comme $e^{x-1} > 0$ sur $\R$, le signe de $\varphi'(x)$ est celui du polynôme du second degré $x(x+2)$, dont les racines sont $x = -2$ et $x = 0$ :
        \begin{itemize}
            \item Sur $]-\infty, -2[$, $\varphi'(x) > 0$, donc $\varphi$ est strictement croissante.
            \item Sur $[-2, 0]$, $\varphi'(x) \le 0$, donc $\varphi$ est strictement décroissante.
            \item Sur $[0, +\infty[$, $\varphi'(x) \ge 0$, donc $\varphi$ est strictement croissante.
        \end{itemize}
        Étudions maintenant les valeurs prises par $\varphi$ :
        \begin{itemize}
            \item \textbf{Cas $x \le 0$ :}
            $\varphi(0) = 0^2 e^{-1} = 0$.
            En $x = -2$, $\varphi$ atteint son maximum local sur $]-\infty, 0]$ :
            \[
            \varphi(-2) = (-2)^2 e^{-2-1} = 4e^{-3} = \frac{4}{e^3}.
            \]
            Or $e > 2{,}7 \implies e^3 > 2{,}7^3 = 19{,}683 > 4$. Ainsi :
            \[
            \varphi(-2) = \frac{4}{e^3} < \frac{4}{19{,}683} \approx 0{,}199 < 1.
            \]
            Par conséquent, pour tout $x \le 0$, on a $\varphi(x) \le \frac{4}{e^3} < 1$. L'équation $\varphi(x) = 1$ n'admet donc \textbf{aucune solution négative ou nulle}.
            \item \textbf{Cas $x > 0$ :}
            Sur $[0, +\infty[$, $\varphi$ est continue et \textbf{strictement croissante}.
            Elle réalise une bijection de $[0, +\infty[$ sur $[\varphi(0), \lim_{x \to +\infty} \varphi(x)[ = [0, +\infty[$.
            Comme $1 \in [0, +\infty[$, d'après le corollaire du théorème des valeurs intermédiaires (théorème de la bijection), l'équation $\varphi(x) = 1$ admet une \textbf{unique solution} sur $[0, +\infty[$.
        \end{itemize}
        Or nous avons déjà vérifié que $\varphi(1) = 1$.
        L'unique solution réelle de l'équation $a^2 e^{a-1} = 1$ est donc $\boxed{a = 1}$.
    \end{itemize}

    \item \textbf{Équation de la tangente et représentation graphique :}
    Pour $a = 1$, la pente vaut :
    \[
    f'(1) = (1+1)e^{1-1} = 2 \times 1 = 2.
    \]
    Le point de contact est $T(1, f(1))$ avec $f(1) = 1 \cdot e^0 + 1 = 2$.
    L'ordonnée à l'origine étant nulle par construction, l'équation de la tangente est :
    \[
    \boxed{y = 2x}.
    \]
\end{enumerate}
\end{correctionbox}

\begin{center}
\begin{tikzpicture}[scale=1.1]
  \begin{axis}[
    axis lines=middle,
    xlabel={$x$},
    ylabel={$y$},
    xmin=-2.8, xmax=2.5,
    ymin=-0.8, ymax=4.8,
    grid=both,
    grid style={line width=.1pt, draw=gray!25},
    major grid style={line width=.2pt, draw=gray!40},
    axis line style={-Latex[round], thick},
    legend pos=north west,
    legend cell align={left},
    samples=150,
    width=13cm, height=8.5cm
  ]
    % Courbe de f(x) = x * exp(x-1) + 1
    \addplot[domain=-2.5:1.6, color=llgblue, very thick, smooth] {x*exp(x-1) + 1};
    \addlegendentry{$\mathcal{C}_f : y = x e^{x-1} + 1$};

    % Tangente en a = 1 : y = 2x
    \addplot[domain=-0.4:2.3, color=red!80!black, thick, dashed] {2*x};
    \addlegendentry{Tangente $T_1 : y = 2x$};

    % Points remarquables
    \node[circle, fill=red, inner sep=1.8pt, label={below right:\small $T(1, 2)$}] at (axis cs:1, 2) {};
    \node[circle, fill=black, inner sep=1.5pt, label={below left:\small $O(0, 0)$}] at (axis cs:0, 0) {};
    \node[circle, fill=blue!70!black, inner sep=1.5pt, label={above left:\small $(0, 1)$}] at (axis cs:0, 1) {};
  \end{axis}
\end{tikzpicture}
\end{center}

% =========================================================================
% EXERCICE 6
% =========================================================================
\subsection*{Exercice 6 : Tangente commune à deux courbes exponentielles}
\addcontentsline{toc}{subsection}{Exercice 6 : Tangente commune à deux courbes exponentielles}

On considère les deux fonctions $f$ et $g$ définies sur $\R$ par :
\[
f(x) = e^x \qquad \text{et} \qquad g(x) = 1 - e^{-x}.
\]
On note respectivement $\mathcal{C}_f$ et $\mathcal{C}_g$ leurs courbes représentatives dans un repère orthonormé $(O; \vec{i}, \vec{j})$.
\begin{enumerate}[label=\textbf{\arabic*.}]
    \item Soit $a \in \R$. Déterminer l'équation de la tangente $T_a$ à la courbe $\mathcal{C}_f$ au point d'abscisse $a$.
    \item Soit $b \in \R$. Déterminer l'équation de la tangente $D_b$ à la courbe $\mathcal{C}_g$ au point d'abscisse $b$.
    \item Montrer que les droites $T_a$ et $D_b$ sont confondues si et seulement si $b = -a$ et $a$ vérifie la relation :
    \[
    2(a-1)e^a + 1 = 0.
    \]
    \item Étudier les variations de la fonction $h : x \mapsto 2(x-1)e^x + 1$ sur $\R$ et en déduire que les courbes $\mathcal{C}_f$ et $\mathcal{C}_g$ admettent \textbf{exactement deux tangentes communes distinctes}.
    \item Mettre en évidence la symétrie centrale liant les deux courbes et interpréter géométriquement la configuration.
\end{enumerate}

\begin{correctionbox}[Exercice 6]
\begin{enumerate}[label=\textbf{\arabic*.}]
    \item \textbf{Équation de la tangente $T_a$ à $\mathcal{C}_f$ :}
    $f$ est dérivable sur $\R$ avec $f'(x) = e^x$.
    Au point d'abscisse $a$, le coefficient directeur est $f'(a) = e^a$ et l'ordonnée est $f(a) = e^a$.
    \[
    T_a : y = f'(a)(x - a) + f(a) = e^a(x - a) + e^a = e^a x - a e^a + e^a.
    \]
    D'où :
    \[
    \boxed{T_a : y = e^a x + (1 - a)e^a}.
    \]

    \item \textbf{Équation de la tangente $D_b$ à $\mathcal{C}_g$ :}
    $g$ est dérivable sur $\R$. Sa dérivée vaut :
    \[
    g'(x) = 0 - (-e^{-x}) = e^{-x}.
    \]
    Au point d'abscisse $b$, on a $g(b) = 1 - e^{-b}$ et $g'(b) = e^{-b}$. L'équation de la tangente $D_b$ est :
    \[
    D_b : y = g'(b)(x - b) + g(b) = e^{-b}(x - b) + 1 - e^{-b} = e^{-b} x - b e^{-b} + 1 - e^{-b}.
    \]
    D'où :
    \[
    \boxed{D_b : y = e^{-b} x + \bigl(1 - (b+1)e^{-b}\bigr)}.
    \]

    \item \textbf{Condition d'identification des deux tangentes :}
    Deux droites non verticales sont confondues si et seulement si elles ont le même coefficient directeur et la même ordonnée à l'origine. Ainsi, $T_a = D_b$ équivaut au système :
    \[
    \begin{cases}
    e^a = e^{-b} & \text{(Égalité des pentes)} \\
    (1 - a)e^a = 1 - (b+1)e^{-b} & \text{(Égalité des ordonnées à l'origine)}
    \end{cases}
    \]
    \begin{itemize}
        \item De la première équation, par injectivité de la fonction exponentielle :
        \[
        e^a = e^{-b} \iff a = -b \iff b = -a.
        \]
        \item Substituons $b = -a$ et $e^{-b} = e^a$ dans la seconde équation :
        \[
        (1 - a)e^a = 1 - (-a + 1)e^a \iff (1 - a)e^a = 1 - (1 - a)e^a.
        \]
        En regroupant les termes en $e^a$ dans le membre de gauche :
        \[
        2(1 - a)e^a = 1 \iff 2(1 - a)e^a - 1 = 0 \iff -2(a - 1)e^a - 1 = 0.
        \]
        En multipliant par $-1$, on obtient exactement l'équation requise :
        \[
        \boxed{2(a - 1)e^a + 1 = 0}.
        \]
    \end{itemize}

    \item \textbf{Étude de la fonction $h(x) = 2(x-1)e^x + 1$ et nombre de tangentes communes :}
    \begin{itemize}
        \item \textbf{Dérivabilité et calcul de dérivée :}
        $h$ est dérivable sur $\R$ comme produit et somme de fonctions usuelles :
        \[
        h'(x) = 2(1)e^x + 2(x-1)e^x + 0 = 2e^x + 2xe^x - 2e^x = 2x e^x.
        \]
        \item \textbf{Signe de $h'(x)$ :}
        Puisque pour tout $x \in \R$, $2e^x > 0$, le signe de $h'(x)$ est strictement identique à celui de $x$ :
        \[
        h'(x) < 0 \text{ pour } x < 0, \qquad h'(0) = 0, \qquad h'(x) > 0 \text{ pour } x > 0.
        \]
        \item \textbf{Variations et extremum :}
        $h$ est strictement décroissante sur $]-\infty, 0]$ et strictement croissante sur $[0, +\infty[$.
        Elle admet en $x = 0$ un minimum absolu qui vaut :
        \[
        h(0) = 2(0 - 1)e^0 + 1 = -2 + 1 = -1 < 0.
        \]
        \item \textbf{Limites aux bornes :}
        \begin{itemize}
            \item En $-\infty$ : Développons $h(x) = 2x e^x - 2e^x + 1$.
            D'après les croissances comparées, $\lim_{x \to -\infty} x e^x = 0$ et $\lim_{x \to -\infty} e^x = 0$.
            Donc $\lim_{x \to -\infty} h(x) = 0 - 0 + 1 = 1 > 0$.
            \item En $+\infty$ : $\lim_{x \to +\infty} 2(x-1) = +\infty$ et $\lim_{x \to +\infty} e^x = +\infty$, donc par produit $\lim_{x \to +\infty} h(x) = +\infty$.
        \end{itemize}
        \item \textbf{Application du corollaire du TVI (Théorème de la bijection) :}
        \begin{itemize}
            \item Sur $]-\infty, 0]$, $h$ est continue et strictement décroissante.
            Elle réalise une bijection de $]-\infty, 0]$ sur $[h(0), \lim_{-\infty} h[ = [-1, 1[$.
            Comme $0 \in ]-1, 1[$, il existe un \textbf{unique réel $a_1 \in ]-\infty, 0[$} tel que $h(a_1) = 0$.
            Encadrement : $h(-1) = 2(-2)e^{-1} + 1 = 1 - 4/e \approx -0{,}47 < 0$, et $h(-2) = 2(-3)e^{-2} + 1 = 1 - 6/e^2 \approx 1 - 0{,}81 = +0{,}19 > 0$.
            Donc $-2 < a_1 < -1$ (numériquement, $a_1 \approx -1{,}678$).
            \item Sur $[0, +\infty[$, $h$ est continue et strictement croissante.
            Elle réalise une bijection de $[0, +\infty[$ sur $[-1, +\infty[$.
            Comme $0 \in ]-1, +\infty[$, il existe un \textbf{unique réel $a_2 \in ]0, +\infty[$} tel que $h(a_2) = 0$.
            Encadrement : $h(0) = -1 < 0$ et $h(1) = 2(0)e^1 + 1 = 1 > 0$.
            Donc $0 < a_2 < 1$ (numériquement, $a_2 \approx 0{,}768$).
        \end{itemize}
        \item \textbf{Conclusion :}
        L'équation admet exactement \textbf{deux solutions réelles distinctes} $a_1$ et $a_2$.
        À chacune de ces valeurs correspond une abscisse de contact $a$ sur $\mathcal{C}_f$ et une abscisse de contact $b = -a$ sur $\mathcal{C}_g$.
        Les courbes $\mathcal{C}_f$ et $\mathcal{C}_g$ admettent donc \textbf{exactement deux tangentes communes distinctes}.
    \end{itemize}

    \item \textbf{Interprétation géométrique et centre de symétrie :}
    Considérons la transformation géométrique de centre $\Omega\left(0, \frac{1}{2}\right)$.
    Pour tout $x \in \R$ :
    \[
    \frac{f(x) + g(-x)}{2} = \frac{e^x + (1 - e^{-(-x)})}{2} = \frac{e^x + 1 - e^x}{2} = \frac{1}{2}.
    \]
    La courbe $\mathcal{C}_g$ est donc l'\textbf{image de la courbe $\mathcal{C}_f$ par la symétrie centrale} de centre $\Omega\left(0, \frac{1}{2}\right)$.
    Puisque la symétrie centrale conserve l'alignement et les contacts de tangence, toute tangente commune à $\mathcal{C}_f$ et $\mathcal{C}_g$ est soit globalement invariante par cette symétrie (passant par $\Omega$), soit échangée avec l'autre tangente commune. Les deux tangentes communes sont en fait symétriques l'une de l'autre par rapport à $\Omega$.
\end{enumerate}
\end{correctionbox}

\begin{center}
\begin{tikzpicture}[scale=1.1]
  \begin{axis}[
    axis lines=middle,
    xlabel={$x$},
    ylabel={$y$},
    xmin=-3.2, xmax=3.2,
    ymin=-2.5, ymax=4.5,
    grid=both,
    grid style={line width=.1pt, draw=gray!25},
    major grid style={line width=.2pt, draw=gray!40},
    axis line style={-Latex[round], thick},
    legend pos=north west,
    legend cell align={left},
    samples=150,
    width=13cm, height=8.5cm
  ]
    % f(x) = exp(x)
    \addplot[domain=-3.2:1.5, color=llgblue, very thick, smooth] {exp(x)};
    \addlegendentry{$\mathcal{C}_f : y = e^x$};

    % g(x) = 1 - exp(-x)
    \addplot[domain=-1.5:3.2, color=emeraldframe, very thick, smooth] {1 - exp(-x)};
    \addlegendentry{$\mathcal{C}_g : y = 1 - e^{-x}$};

    % Tangente 1 pour a2 approx 0.768 : pente exp(0.768) approx 2.155, ordonnee (1-0.768)*2.155 approx 0.500
    \addplot[domain=-1.2:1.8, color=red!80!black, thick, dashed] {2.155*x + 0.500};
    \addlegendentry{Tangente commune 1 ($a_2 \approx 0{,}77$)};

    % Tangente 2 pour a1 approx -1.678 : pente exp(-1.678) approx 0.187, ordonnee (1 - (-1.678))*0.187 approx 0.500
    \addplot[domain=-3.0:3.0, color=purple!80!black, thick, dotted] {0.187*x + 0.500};
    \addlegendentry{Tangente commune 2 ($a_1 \approx -1{,}68$)};

    % Centre de symétrie Omega(0, 0.5)
    \node[circle, fill=orange!90!black, inner sep=2pt, label={above left:\small $\Omega(0, 0{,}5)$}] at (axis cs:0, 0.5) {};
  \end{axis}
\end{tikzpicture}
\end{center}

% =========================================================================
% EXERCICE 7
% =========================================================================
\subsection*{Exercice 7 : Étude approfondie de $f(x) = \dfrac{e^x - 1}{e^x - x}$ et position relative}
\addcontentsline{toc}{subsection}{Exercice 7 : Étude de $f(x) = \dfrac{e^x - 1}{e^x - x}$}

On s'intéresse à la fonction $f$ définie formellement par :
\[
f(x) = \frac{e^x - 1}{e^x - x}.
\]
\begin{enumerate}[label=\textbf{\arabic*.}]
    \item Démontrer que pour tout réel $x \in \R$, $e^x > x$ en distinguant soigneusement les cas $x \le 0$ et $x > 0$. En déduire l'ensemble de définition $\mathcal{D}_f$.
    \item Étudier les variations de la fonction auxiliaire $g(x) = e^x - x - 1$ sur $\R$ et en déduire son signe.
    \item Étudier la position relative de la courbe représentative $\mathcal{C}_f$ et de la droite $\Delta$ d'équation $y = x$. On établira au préalable que pour tout $x \in \R$ :
    \[
    f(x) - x = \frac{(1 - x)g(x)}{e^x - x}.
    \]
    \item Préciser les points d'intersection de $\mathcal{C}_f$ et $\Delta$ ainsi que l'allure géométrique de la courbe au voisinage de ces points.
\end{enumerate}

\begin{correctionbox}[Exercice 7]
\begin{enumerate}[label=\textbf{\arabic*.}]
    \item \textbf{Démonstration de $e^x > x$ sur $\R$ et ensemble de définition :}
    \begin{itemize}
        \item \textbf{Premier cas : $x \le 0$ :}
        D'après les propriétés du cours, la fonction exponentielle est strictement positive sur $\R$, donc $e^x > 0$.
        Or pour tout $x \le 0$, on a par transitivité :
        \[
        e^x > 0 \ge x \implies e^x > x.
        \]
        \item \textbf{Second cas : $x > 0$ :}
        Considérons la fonction d'écart $d(x) = e^x - x$ sur l'intervalle $[0, +\infty[$.
        $d$ est dérivable sur $[0, +\infty[$ de dérivée $d'(x) = e^x - 1$.
        Pour tout $x > 0$, par stricte croissance de l'exponentielle, $e^x > e^0 = 1$, d'où $d'(x) > 0$.
        La fonction $d$ est donc \textbf{strictement croissante} sur $[0, +\infty[$.
        Par conséquent, pour tout $x > 0$ :
        \[
        d(x) > d(0) = e^0 - 0 = 1 > 0 \implies e^x - x > 0 \iff e^x > x.
        \]
        \item \textbf{Bilan :} L'inégalité $e^x > x$ est vraie pour tout $x \in \R$.
        On en déduit que pour tout $x \in \R$, $e^x - x \neq 0$. Le dénominateur ne s'annule donc jamais sur $\R$.
        L'ensemble de définition de la fonction $f$ est ainsi :
        \[
        \boxed{\mathcal{D}_f = \R}.
        \]
    \end{itemize}

    \item \textbf{Étude de la fonction $g(x) = e^x - x - 1$ :}
    \begin{itemize}
        \item $g$ est dérivable sur $\R$ comme somme de fonctions dérivables, et :
        \[
        g'(x) = e^x - 1.
        \]
        \item Étude du signe de $g'(x)$ :
        \[
        g'(x) = 0 \iff e^x = 1 \iff x = 0,
        \]
        \[
        g'(x) > 0 \iff e^x > 1 \iff x > 0,
        \]
        \[
        g'(x) < 0 \iff e^x < 1 \iff x < 0.
        \]
        \item $g$ est donc strictement décroissante sur $]-\infty, 0]$ et strictement croissante sur $[0, +\infty[$.
        Elle admet en $x = 0$ un \textbf{minimum absolu} valant :
        \[
        g(0) = e^0 - 0 - 1 = 1 - 1 = 0.
        \]
        \item \textbf{Conclusion pour le signe de $g$ :}
        Pour tout $x \in \R$, $g(x) \ge 0$, l'égalité $g(x) = 0$ n'ayant lieu qu'en $x = 0$ :
        \[
        \forall x \in \R \setminus \{0\}, \quad g(x) > 0 \qquad \text{et} \qquad g(0) = 0.
        \]
    \end{itemize}

    \item \textbf{Position relative de $\mathcal{C}_f$ et de la droite $\Delta : y = x$ :}
    \begin{itemize}
        \item \textbf{Établissement de l'identité algébrique :}
        Calculons la différence $f(x) - x$ en réduisant au même dénominateur :
        \[
        f(x) - x = \frac{e^x - 1}{e^x - x} - x = \frac{e^x - 1 - x(e^x - x)}{e^x - x} = \frac{e^x - 1 - x e^x + x^2}{e^x - x}.
        \]
        Développons d'autre part le produit proposé au numérateur $(1 - x)g(x)$ :
        \[
        (1 - x)g(x) = (1 - x)(e^x - x - 1) = 1 \cdot (e^x - x - 1) - x(e^x - x - 1)
        \]
        \[
        = e^x - x - 1 - x e^x + x^2 + x = e^x - 1 - x e^x + x^2.
        \]
        Les deux numérateurs sont identiques, d'où pour tout $x \in \R$ :
        \[
        \boxed{f(x) - x = \frac{(1 - x)g(x)}{e^x - x}}.
        \]
        \item \textbf{Étude du signe du quotient :}
        \begin{itemize}
            \item Le dénominateur $e^x - x$ est strictement positif sur $\R$ (question 1).
            \item Le facteur $g(x)$ est strictement positif pour tout $x \neq 0$, et s'annule en $x = 0$ (question 2).
            \item Par conséquent, pour tout $x \neq 0$, le signe de $f(x) - x$ est \textbf{strictement identique à celui de $(1 - x)$}.
        \end{itemize}
        Dressons le tableau récapitulatif des signes :
        \begin{center}
        \begin{tabular}{|c|lccccccr|}
        \hline
        $x$ & $-\infty$ & & $0$ & & $1$ & & $+\infty$ \\
        \hline
        $1 - x$ & & $+$ & \vline & $+$ & $0$ & $-$ & \\
        \hline
        $g(x)$ & & $+$ & $0$ & $+$ & \vline & $+$ & \\
        \hline
        $e^x - x$ & & $+$ & \vline & $+$ & \vline & $+$ & \\
        \hline
        $f(x) - x$ & & $+$ & $0$ & $+$ & $0$ & $-$ & \\
        \hline
        \end{tabular}
        \end{center}
        \item \textbf{Conclusions géométriques :}
        \begin{itemize}
            \item Sur $]-\infty, 0[ \cup ]0, 1[$, $f(x) - x > 0$ : la courbe $\mathcal{C}_f$ est \textbf{strictement au-dessus} de la droite $\Delta$.
            \item Sur $]1, +\infty[$, $f(x) - x < 0$ : la courbe $\mathcal{C}_f$ est \textbf{strictement au-dessous} de la droite $\Delta$.
            \item En $x = 0$, $f(0) - 0 = 0$ : $\mathcal{C}_f$ et $\Delta$ se coupent à l'origine $O(0, 0)$.
            \item En $x = 1$, $f(1) - 1 = \dfrac{e - 1}{e - 1} - 1 = 0$ : $\mathcal{C}_f$ et $\Delta$ se coupent au point $A(1, 1)$.
        \end{itemize}
    \end{itemize}

    \item \textbf{Comportement aux points d'intersection :}
    \begin{itemize}
        \item En $x = 0$, $f(x) - x$ s'annule sans changer de signe (il reste strictement positif de part et d'autre).
        La droite $\Delta : y = x$ est donc \textbf{tangente} à la courbe $\mathcal{C}_f$ en l'origine $O(0, 0)$, sans la traverser.
        Vérification par la dérivée : $f'(0) = 1$.
        \item En $x = 1$, $f(x) - x$ s'annule en changeant de signe (de positif à négatif).
        La courbe $\mathcal{C}_f$ \textbf{traverse} la droite $\Delta$ au point $A(1, 1)$.
    \end{itemize}
\end{enumerate}
\end{correctionbox}

\begin{center}
\begin{tikzpicture}[scale=1.1]
  \begin{axis}[
    axis lines=middle,
    xlabel={$x$},
    ylabel={$y$},
    xmin=-3.2, xmax=3.5,
    ymin=-2.2, ymax=3.2,
    grid=both,
    grid style={line width=.1pt, draw=gray!25},
    major grid style={line width=.2pt, draw=gray!40},
    axis line style={-Latex[round], thick},
    legend pos=north west,
    legend cell align={left},
    samples=150,
    width=13cm, height=8.5cm
  ]
    % Courbe f(x) = (exp(x)-1)/(exp(x)-x)
    \addplot[domain=-3:3.2, color=llgblue, very thick, smooth] {(exp(x)-1)/(exp(x)-x)};
    \addlegendentry{$\mathcal{C}_f : y = \dfrac{e^x-1}{e^x-x}$};

    % Première bissectrice y = x
    \addplot[domain=-2.2:3.0, color=red!80!black, thick, dashed] {x};
    \addlegendentry{Droite $\Delta : y = x$};

    % Asymptotes horizontales : y = 1 en +infini et y = 0 en -infini
    \addplot[domain=-3.2:3.5, color=gray!80!black, dotted, thick] {1};
    \addlegendentry{Asymptote $y = 1$};

    % Points remarquables
    \node[circle, fill=black, inner sep=1.8pt, label={below right:\small $O(0, 0)$ (tangence)}] at (axis cs:0, 0) {};
    \node[circle, fill=blue!70!black, inner sep=1.8pt, label={above left:\small $A(1, 1)$ (traversée)}] at (axis cs:1, 1) {};
  \end{axis}
\end{tikzpicture}
\end{center}

% =========================================================================
% EXERCICE 8
% =========================================================================
\subsection*{Exercice 8 : Famille de fonctions exponentielles $f_k(x) = (x+1)e^{kx}$ ($k \in \Z$)}
\addcontentsline{toc}{subsection}{Exercice 8 : Famille de fonctions $f_k(x) = (x+1)e^{kx}$}

Pour tout entier relatif $k \in \Z$, on considère la fonction $f_k$ définie sur $\R$ par :
\[
f_k(x) = (x + 1)e^{kx}.
\]
On note $\mathcal{C}_k$ sa courbe représentative dans un repère orthonormé $(O; \vec{i}, \vec{j})$.
\begin{enumerate}[label=\textbf{\arabic*.}]
    \item Démontrer que toutes les courbes $\mathcal{C}_k$ passent par exactement deux points fixes dont on précisera les coordonnées.
    \item Pour tout $k \in \Z$, étudier la position relative de la courbe $\mathcal{C}_{k+1}$ par rapport à la courbe $\mathcal{C}_k$.
    \item Étudier les variations de la fonction $f_k$ en discutant selon le signe de l'entier $k$.
    \item Déterminer les coordonnées de l'extremum éventuel de $f_k$ et étudier le lieu géométrique de ces extrema lorsque $k$ varie.
\end{enumerate}

\begin{correctionbox}[Exercice 8]
\begin{enumerate}[label=\textbf{\arabic*.}]
    \item \textbf{Recherche des points fixes universels :}
    Un point $M(x, y)$ appartient à toutes les courbes $\mathcal{C}_k$ si et seulement si pour tous entiers distincts $k, k' \in \Z$, on a :
    \[
    f_k(x) = f_{k'}(x) \iff (x + 1)e^{kx} = (x + 1)e^{k'x} \iff (x + 1)\bigl(e^{kx} - e^{k'x}\bigr) = 0.
    \]
    Cette équation produit est vérifiée pour tous $k \neq k'$ si et seulement si l'une des deux conditions suivantes est satisfaite :
    \begin{itemize}
        \item \textbf{Premier facteur nul :} $x + 1 = 0 \iff x = -1$.
        Pour $x = -1$, on a $f_k(-1) = (-1 + 1)e^{-k} = 0$, et ce pour \textbf{tout entier $k \in \Z$}.
        Toutes les courbes passent donc par le point fixe $\boxed{A(-1, 0)}$.
        \item \textbf{Second facteur nul pour tous $k, k'$ :} $e^{kx} = e^{k'x} \iff e^{(k - k')x} = 1 \iff (k - k')x = 0$.
        Puisque cette égalité doit avoir lieu pour $k \neq k'$, on en déduit nécessairement $x = 0$.
        Pour $x = 0$, on a $f_k(0) = (0 + 1)e^0 = 1$, et ce pour \textbf{tout entier $k \in \Z$}.
        Toutes les courbes passent donc par le point fixe $\boxed{B(0, 1)}$.
    \end{itemize}
    \textbf{Conclusion :} L'ensemble des courbes $\mathcal{C}_k$ admet exactement \textbf{deux points fixes communs} : $A(-1, 0)$ et $B(0, 1)$.

    \item \textbf{Position relative de $\mathcal{C}_{k+1}$ par rapport à $\mathcal{C}_k$ :}
    Considérons la différence $d_k(x) = f_{k+1}(x) - f_k(x)$ pour tout $x \in \R$ :
    \[
    d_k(x) = (x+1)e^{(k+1)x} - (x+1)e^{kx} = (x+1)e^{kx}e^x - (x+1)e^{kx} = (x+1)e^{kx}(e^x - 1).
    \]
    Comme pour tout $x \in \R$, $e^{kx} > 0$, le signe de $d_k(x)$ est le produit des signes de $(x+1)$ et $(e^x - 1)$ :
    \begin{itemize}
        \item Le facteur $(x+1)$ s'annule en $-1$, est négatif pour $x < -1$ et positif pour $x > -1$.
        \item Le facteur $(e^x - 1)$ s'annule en $0$, est négatif pour $x < 0$ et positif pour $x > 0$.
    \end{itemize}
    Dressons le tableau de signes de $d_k(x)$ :
    \begin{center}
    \begin{tabular}{|c|lccccccr|}
    \hline
    $x$ & $-\infty$ & & $-1$ & & $0$ & & $+\infty$ \\
    \hline
    $x + 1$ & & $-$ & $0$ & $+$ & \vline & $+$ & \\
    \hline
    $e^x - 1$ & & $-$ & \vline & $-$ & $0$ & $+$ & \\
    \hline
    $e^{kx}$ & & $+$ & \vline & $+$ & \vline & $+$ & \\
    \hline
    $d_k(x)$ & & $+$ & $0$ & $-$ & $0$ & $+$ & \\
    \hline
    \end{tabular}
    \end{center}
    On en déduit la position relative universelle (valable pour tout $k \in \Z$) :
    \begin{itemize}
        \item Sur $]-\infty, -1[$, $d_k(x) > 0$ : $\mathcal{C}_{k+1}$ est \textbf{strictement au-dessus} de $\mathcal{C}_k$.
        \item En $x = -1$, $d_k(-1) = 0$ : intersection au point fixe $A(-1, 0)$.
        \item Sur $]-1, 0[$, $d_k(x) < 0$ : $\mathcal{C}_{k+1}$ est \textbf{strictement au-dessous} de $\mathcal{C}_k$.
        \item En $x = 0$, $d_k(0) = 0$ : intersection au point fixe $B(0, 1)$.
        \item Sur $]0, +\infty[$, $d_k(x) > 0$ : $\mathcal{C}_{k+1}$ est \textbf{strictement au-dessus} de $\mathcal{C}_k$.
    \end{itemize}

    \item \textbf{Variations de $f_k$ selon le signe de $k \in \Z$ :}
    $f_k$ est dérivable sur $\R$ comme produit de fonctions dérivables. Par $(uv)' = u'v + uv'$ :
    \[
    f_k'(x) = 1 \cdot e^{kx} + (x+1) \cdot k e^{kx} = \bigl[1 + k(x+1)\bigr]e^{kx} = (kx + k + 1)e^{kx}.
    \]
    Le facteur $e^{kx}$ étant strictement positif, le signe de $f_k'(x)$ est celui de l'expression affine $kx + k + 1$.
    \begin{itemize}
        \item \textbf{Cas 1 : $k = 0$ :}
        $f_0(x) = x + 1$. La fonction est affine de pente $1 > 0$, donc \textbf{strictement croissante} sur $\R$.
        \item \textbf{Cas 2 : $k > 0$ :}
        Le coefficient directeur $k$ est strictement positif.
        \[
        kx + k + 1 = 0 \iff kx = -(k+1) \iff x = -\frac{k+1}{k} = -1 - \frac{1}{k}.
        \]
        Puisque $k > 0$ :
        \begin{itemize}
            \item Pour $x < -1 - \frac{1}{k}$, $f_k'(x) < 0$, donc $f_k$ est strictement décroissante.
            \item Pour $x > -1 - \frac{1}{k}$, $f_k'(x) > 0$, donc $f_k$ est strictement croissante.
        \end{itemize}
        $f_k$ admet donc un \textbf{minimum absolu} en $x_k = -1 - \frac{1}{k}$.
        \item \textbf{Cas 3 : $k < 0$ :}
        Posons $k = -m$ avec $m \in \N^*$. Le coefficient directeur $k < 0$ est strictement négatif.
        L'annulation a toujours lieu en $x_k = -1 - \frac{1}{k} = -1 + \frac{1}{m} > -1$.
        \begin{itemize}
            \item Pour $x < -1 - \frac{1}{k}$, $f_k'(x) > 0$, donc $f_k$ est strictement croissante.
            \item Pour $x > -1 - \frac{1}{k}$, $f_k'(x) < 0$, donc $f_k$ est strictement décroissante.
        \end{itemize}
        $f_k$ admet donc un \textbf{maximum absolu} en $x_k = -1 - \frac{1}{k}$.
    \end{itemize}

    \item \textbf{Coordonnées des extrema et comportement asymptotique :}
    Calculons la valeur de l'extremum en $x_k = -1 - \dfrac{1}{k}$ :
    \[
    f_k(x_k) = \left(-1 - \frac{1}{k} + 1\right)e^{k\left(-1 - \frac{1}{k}\right)} = -\frac{1}{k} e^{-k - 1} = -\frac{1}{k e^{k+1}}.
    \]
    Notons $S_k\left(-1 - \frac{1}{k}, -\frac{1}{k e^{k+1}}\right)$ le sommet de la courbe $\mathcal{C}_k$.
    \begin{itemize}
        \item Si $k \to +\infty$ : $x_k = -1 - \frac{1}{k} \to -1$ et $y_k = -\frac{1}{k e^{k+1}} \to 0$.
        \item Si $k \to -\infty$ : $x_k = -1 - \frac{1}{k} \to -1$ et $y_k = -\frac{1}{k e^{k+1}} \to 0$ (car $k e^{k+1} \to -\infty$).
    \end{itemize}
    Lorsque $|k| \to +\infty$, les sommets des courbes $\mathcal{C}_k$ \textbf{convergent vers le point fixe $A(-1, 0)$}.
\end{enumerate}
\end{correctionbox}

\begin{center}
\begin{tikzpicture}[scale=1.1]
  \begin{axis}[
    axis lines=middle,
    xlabel={$x$},
    ylabel={$y$},
    xmin=-2.6, xmax=2.2,
    ymin=-1.5, ymax=4.2,
    grid=both,
    grid style={line width=.1pt, draw=gray!25},
    major grid style={line width=.2pt, draw=gray!40},
    axis line style={-Latex[round], thick},
    legend pos=north west,
    legend cell align={left},
    samples=150,
    width=13cm, height=8.5cm
  ]
    % Courbe k = -2 : f_{-2}(x) = (x+1)exp(-2x)
    \addplot[domain=-1.2:2.2, color=purple!80!black, thick, smooth] {(x+1)*exp(-2*x)};
    \addlegendentry{$k = -2$};

    % Courbe k = -1 : f_{-1}(x) = (x+1)exp(-x)
    \addplot[domain=-1.5:2.2, color=blue!70!black, thick, smooth] {(x+1)*exp(-x)};
    \addlegendentry{$k = -1$};

    % Courbe k = 0 : f_0(x) = x+1
    \addplot[domain=-2.4:2.2, color=gray!80!black, thick, dashed] {x+1};
    \addlegendentry{$k = 0$};

    % Courbe k = 1 : f_1(x) = (x+1)exp(x)
    \addplot[domain=-2.5:1.2, color=emeraldframe, thick, smooth] {(x+1)*exp(x)};
    \addlegendentry{$k = 1$};

    % Courbe k = 2 : f_2(x) = (x+1)exp(2x)
    \addplot[domain=-2.5:0.7, color=red!80!black, thick, smooth] {(x+1)*exp(2*x)};
    \addlegendentry{$k = 2$};

    % Points fixes universels A(-1, 0) et B(0, 1)
    \node[circle, fill=black, inner sep=2.2pt, label={below left:\small $A(-1, 0)$}] at (axis cs:-1, 0) {};
    \node[circle, fill=black, inner sep=2.2pt, label={above left:\small $B(0, 1)$}] at (axis cs:0, 1) {};
  \end{axis}
\end{tikzpicture}
\end{center}

\vspace{0.8cm}
\hrule
\vspace{0.8cm}
\section{Section 4 : Limites et Théorème des croissances comparées}

L'évaluation asymptotique des expressions exponentielles est au cœur de l'analyse réelle. Face aux formes indéterminées (« $\infty - \infty$ », « $\frac{\infty}{\infty}$ », « $0 \times \infty$ »), le réflexe fondamental consiste à isoler le \textbf{terme dominant} (ou prépondérant) et à mobiliser le théorème des croissances comparées, qui formalise le fait que l'exponentielle croît infiniment plus vite que tout polynôme au voisinage de $+\infty$.

% =========================================================================
% EXERCICE 9
% =========================================================================
\subsection*{Exercice 9 : Limites directes en $-\infty$ et $+\infty$ de sept fonctions}
\addcontentsline{toc}{subsection}{Exercice 9 : Limites directes de sept fonctions}

Déterminer la limite en $-\infty$ et en $+\infty$ de chacune des sept fonctions suivantes, et préciser les asymptotes horizontales ou obliques éventuelles :
\begin{enumerate}[label=\textbf{\arabic*.}]
    \item $g_1(x) = e^{-3x}$
    \item $g_2(x) = e^x + e^{-x}$
    \item $g_3(x) = x + e^x$
    \item $g_4(x) = e^{2x} + e^x + 1$
    \item $g_5(x) = e^{3x} - e^x$
    \item $g_6(x) = \dfrac{e^x + 1}{e^x + 2}$
    \item $g_7(x) = \dfrac{-2e^x}{1 + e^x}$
\end{enumerate}

\begin{correctionbox}[Exercice 9]
\begin{enumerate}[label=\textbf{\arabic*.}]
    \item \textbf{Limites de $g_1(x) = e^{-3x}$ :}
    \begin{itemize}
        \item En $-\infty$ : Lorsque $x \to -\infty$, on a $-3x \to +\infty$. Par composition avec $\lim_{X \to +\infty} e^X = +\infty$ :
        \[
        \boxed{\lim_{x \to -\infty} g_1(x) = +\infty}.
        \]
        \item En $+\infty$ : Lorsque $x \to +\infty$, on a $-3x \to -\infty$. Par composition avec $\lim_{X \to -\infty} e^X = 0$ :
        \[
        \boxed{\lim_{x \to +\infty} g_1(x) = 0}.
        \]
        La droite d'équation $y = 0$ (l'axe des abscisses) est \textbf{asymptote horizontale} à la courbe en $+\infty$.
    \end{itemize}

    \item \textbf{Limites de $g_2(x) = e^x + e^{-x}$ :}
    \begin{itemize}
        \item En $-\infty$ : $\lim_{x \to -\infty} e^x = 0$ et $\lim_{x \to -\infty} e^{-x} = +\infty$. Par somme :
        \[
        \boxed{\lim_{x \to -\infty} g_2(x) = +\infty}.
        \]
        \item En $+\infty$ : $\lim_{x \to +\infty} e^x = +\infty$ et $\lim_{x \to +\infty} e^{-x} = 0$. Par somme :
        \[
        \boxed{\lim_{x \to +\infty} g_2(x) = +\infty}.
        \]
        \textit{Remarque :} $g_2(x) - e^x = e^{-x} \to 0$ en $+\infty$, donc la courbe exponentielle $y = e^x$ est asymptote courbe à $\mathcal{C}_{g_2}$ en $+\infty$ (et $y = e^{-x}$ en $-\infty$).
    \end{itemize}

    \item \textbf{Limites de $g_3(x) = x + e^x$ :}
    \begin{itemize}
        \item En $-\infty$ : $\lim_{x \to -\infty} x = -\infty$ et $\lim_{x \to -\infty} e^x = 0$. Par somme de limites :
        \[
        \boxed{\lim_{x \to -\infty} g_3(x) = -\infty}.
        \]
        De plus, $\lim_{x \to -\infty} \bigl(g_3(x) - x\bigr) = \lim_{x \to -\infty} e^x = 0$.
        La droite $\Delta : y = x$ est donc \textbf{asymptote oblique} à la courbe en $-\infty$.
        \item En $+\infty$ : $\lim_{x \to +\infty} x = +\infty$ et $\lim_{x \to +\infty} e^x = +\infty$. Par somme :
        \[
        \boxed{\lim_{x \to +\infty} g_3(x) = +\infty}.
        \]
    \end{itemize}

    \item \textbf{Limites de $g_4(x) = e^{2x} + e^x + 1$ :}
    \begin{itemize}
        \item En $-\infty$ : $\lim_{x \to -\infty} 2x = -\infty \implies \lim_{x \to -\infty} e^{2x} = 0$, et $\lim_{x \to -\infty} e^x = 0$.
        Par somme de limites finies :
        \[
        \boxed{\lim_{x \to -\infty} g_4(x) = 0 + 0 + 1 = 1}.
        \]
        La droite d'équation $y = 1$ est \textbf{asymptote horizontale} en $-\infty$.
        \item En $+\infty$ : $\lim_{x \to +\infty} e^{2x} = +\infty$ et $\lim_{x \to +\infty} e^x = +\infty$. Par somme :
        \[
        \boxed{\lim_{x \to +\infty} g_4(x) = +\infty}.
        \]
    \end{itemize}

    \item \textbf{Limites de $g_5(x) = e^{3x} - e^x$ :}
    \begin{itemize}
        \item En $-\infty$ : $\lim_{x \to -\infty} e^{3x} = 0$ et $\lim_{x \to -\infty} e^x = 0$. Par somme :
        \[
        \boxed{\lim_{x \to -\infty} g_5(x) = 0}.
        \]
        L'axe des abscisses $y = 0$ est \textbf{asymptote horizontale} en $-\infty$.
        \item En $+\infty$ : C'est une forme indéterminée « $\infty - \infty$ ».
        Factorisons par le terme prépondérant $e^{3x}$ :
        \[
        g_5(x) = e^{3x} \left(1 - \frac{e^x}{e^{3x}}\right) = e^{3x}(1 - e^{-2x}).
        \]
        Comme $\lim_{x \to +\infty} (-2x) = -\infty$, $\lim_{x \to +\infty} e^{-2x} = 0$, d'où $\lim_{x \to +\infty} (1 - e^{-2x}) = 1$.
        Comme $\lim_{x \to +\infty} e^{3x} = +\infty$, par produit :
        \[
        \boxed{\lim_{x \to +\infty} g_5(x) = +\infty}.
        \]
    \end{itemize}

    \item \textbf{Limites de $g_6(x) = \dfrac{e^x + 1}{e^x + 2}$ :}
    \begin{itemize}
        \item En $-\infty$ : $\lim_{x \to -\infty} e^x = 0$. Donc par quotient de limites :
        \[
        \boxed{\lim_{x \to -\infty} g_6(x) = \frac{0 + 1}{0 + 2} = \frac{1}{2}}.
        \]
        La droite $y = \frac{1}{2}$ est \textbf{asymptote horizontale} en $-\infty$.
        \item En $+\infty$ : Forme indéterminée « $\frac{\infty}{\infty}$ ».
        Factorisons par le terme dominant $e^x$ au numérateur et au dénominateur :
        \[
        g_6(x) = \frac{e^x\left(1 + \frac{1}{e^x}\right)}{e^x\left(1 + \frac{2}{e^x}\right)} = \frac{1 + e^{-x}}{1 + 2e^{-x}}.
        \]
        Puisque $\lim_{x \to +\infty} e^{-x} = 0$, il vient :
        \[
        \boxed{\lim_{x \to +\infty} g_6(x) = \frac{1 + 0}{1 + 0} = 1}.
        \]
        La droite $y = 1$ est \textbf{asymptote horizontale} en $+\infty$.
    \end{itemize}

    \item \textbf{Limites de $g_7(x) = \dfrac{-2e^x}{1 + e^x}$ :}
    \begin{itemize}
        \item En $-\infty$ : $\lim_{x \to -\infty} e^x = 0$, d'où $\lim_{x \to -\infty} (-2e^x) = 0$ et $\lim_{x \to -\infty} (1 + e^x) = 1$.
        Par quotient :
        \[
        \boxed{\lim_{x \to -\infty} g_7(x) = \frac{0}{1} = 0}.
        \]
        La droite $y = 0$ est \textbf{asymptote horizontale} en $-\infty$.
        \item En $+\infty$ : Forme indéterminée « $\frac{\infty}{\infty}$ ».
        Factorisons par $e^x$ au numérateur et au dénominateur :
        \[
        g_7(x) = \frac{-2e^x}{e^x(e^{-x} + 1)} = \frac{-2}{e^{-x} + 1}.
        \]
        Comme $\lim_{x \to +\infty} e^{-x} = 0$, par quotient de limites finies :
        \[
        \boxed{\lim_{x \to +\infty} g_7(x) = \frac{-2}{0 + 1} = -2}.
        \]
        La droite $y = -2$ est \textbf{asymptote horizontale} en $+\infty$.
    \end{itemize}
\end{enumerate}
\end{correctionbox}

% =========================================================================
% EXERCICE 10
% =========================================================================
\subsection*{Exercice 10 : Limites remarquables et taux de variation}
\addcontentsline{toc}{subsection}{Exercice 10 : Limites remarquables}

Démontrer avec la plus grande rigueur mathématique la valeur des trois limites remarquables suivantes :
\begin{enumerate}[label=\textbf{\arabic*.}]
    \item $\lim_{x \to +\infty} 3x e^{-x} = 0$
    \item $\lim_{x \to -\infty} (x+1)e^x = 0$
    \item $\lim_{x \to 0} \dfrac{1 - e^x}{2x} = -\dfrac{1}{2}$
\end{enumerate}

\begin{correctionbox}[Exercice 10]
\begin{enumerate}[label=\textbf{\arabic*.}]
    \item \textbf{Calcul de $\lim_{x \to +\infty} 3x e^{-x}$ :}
    \begin{itemize}
        \item \textbf{Forme indéterminée initiale :} Lorsque $x \to +\infty$, $3x \to +\infty$ et $e^{-x} \to 0$. C'est une forme indéterminée du type « $\infty \times 0$ ».
        \item \textbf{Transformation algébrique :}
        Puisque $e^{-x} = \dfrac{1}{e^x}$, nous pouvons réécrire l'expression pour tout $x > 0$ sous la forme d'un quotient :
        \[
        3x e^{-x} = 3 \times \frac{x}{e^x} = \frac{3}{\frac{e^x}{x}}.
        \]
        \item \textbf{Application des croissances comparées :}
        D'après le théorème fondamental de croissances comparées en $+\infty$ :
        \[
        \lim_{x \to +\infty} \frac{e^x}{x} = +\infty.
        \]
        Par passage à l'inverse :
        \[
        \lim_{x \to +\infty} \frac{x}{e^x} = 0.
        \]
        Par linéarité de la limite :
        \[
        \boxed{\lim_{x \to +\infty} 3x e^{-x} = 3 \times 0 = 0}.
        \]
    \end{itemize}

    \item \textbf{Calcul de $\lim_{x \to -\infty} (x+1)e^x$ :}
    \begin{itemize}
        \item \textbf{Forme indéterminée initiale :} Lorsque $x \to -\infty$, $x+1 \to -\infty$ et $e^x \to 0$. C'est une forme indéterminée du type « $\infty \times 0$ ».
        \item \textbf{Développement et changement de variable :}
        Développons l'expression :
        \[
        (x+1)e^x = x e^x + e^x.
        \]
        D'une part, par propriété fondamentale de l'exponentielle en $-\infty$ :
        \[
        \lim_{x \to -\infty} e^x = 0.
        \]
        D'autre part, pour évaluer $\lim_{x \to -\infty} x e^x$, effectuons le \textbf{changement de variable} canonique :
        \[
        X = -x \iff x = -X.
        \]
        Lorsque $x \to -\infty$, on a $X \to +\infty$.
        L'expression devient :
        \[
        x e^x = (-X) e^{-X} = -\frac{X}{e^X}.
        \]
        D'après les croissances comparées en $+\infty$, $\lim_{X \to +\infty} \dfrac{X}{e^X} = 0$, donc :
        \[
        \lim_{x \to -\infty} x e^x = \lim_{X \to +\infty} \left(-\frac{X}{e^X}\right) = 0.
        \]
        Par somme de deux limites nulles :
        \[
        \boxed{\lim_{x \to -\infty} (x+1)e^x = 0 + 0 = 0}.
        \]
    \end{itemize}

    \item \textbf{Calcul de $\lim_{x \to 0} \dfrac{1 - e^x}{2x}$ :}
    \begin{itemize}
        \item \textbf{Forme indéterminée initiale :} Lorsque $x \to 0$, le numérateur tend vers $1 - e^0 = 0$ et le dénominateur vers 0. C'est la forme indéterminée « $\frac{0}{0}$ ».
        \item \textbf{Interprétation en terme de taux d'accroissement :}
        Factorisons le coefficient constant $-\dfrac{1}{2}$ :
        \[
        \frac{1 - e^x}{2x} = -\frac{1}{2} \times \frac{e^x - 1}{x} = -\frac{1}{2} \times \frac{e^x - e^0}{x - 0}.
        \]
        On reconnaît exactement le \textbf{taux de variation} en $x_0 = 0$ de la fonction $f : t \mapsto e^t$.
        Puisque la fonction exponentielle est dérivable en $0$, ce taux de variation converge par définition vers le nombre dérivé $f'(0)$ :
        \[
        \lim_{x \to 0} \frac{e^x - e^0}{x - 0} = f'(0).
        \]
        Or $f'(t) = e^t$, d'où $f'(0) = e^0 = 1$.
        Par produit avec le scalaire $-\dfrac{1}{2}$ :
        \[
        \boxed{\lim_{x \to 0} \frac{1 - e^x}{2x} = -\frac{1}{2} \times 1 = -\frac{1}{2}}.
        \]
    \end{itemize}
\end{enumerate}
\end{correctionbox}

% =========================================================================
% EXERCICE 11
% =========================================================================
\subsection*{Exercice 11 : Levée d'indéterminations par factorisation du terme dominant}
\addcontentsline{toc}{subsection}{Exercice 11 : Levée d'indéterminations par terme dominant}

Déterminer la limite en $+\infty$ de chacune des six fonctions suivantes en explicitant la factorisation par le terme prépondérant :
\begin{enumerate}[label=\textbf{\arabic*.}]
    \item $h_1(x) = x^2 + 2 - e^x$
    \item $h_2(x) = \dfrac{2e^x - x}{x^2}$
    \item $h_3(x) = \dfrac{e^x}{x^2 + 1}$
    \item $h_4(x) = e^{2x} - (x+1)e^x$
    \item $h_5(x) = \dfrac{\sqrt{e^x + 1}}{x}$
    \item $h_6(x) = \dfrac{e^{2x} + x^2}{x^2 + x - 3}$
\end{enumerate}

\begin{correctionbox}[Exercice 11]
\begin{enumerate}[label=\textbf{\arabic*.}]
    \item \textbf{Limite de $h_1(x) = x^2 + 2 - e^x$ en $+\infty$ :}
    \begin{itemize}
        \item Forme indéterminée « $\infty - \infty$ ».
        \item Le terme dominant au voisinage de $+\infty$ est l'exponentielle $e^x$. Factorisons par $e^x$ :
        \[
        h_1(x) = e^x \left(\frac{x^2}{e^x} + \frac{2}{e^x} - 1\right).
        \]
        \item Par croissances comparées à l'ordre 2 : $\lim_{x \to +\infty} \dfrac{e^x}{x^2} = +\infty \implies \lim_{x \to +\infty} \dfrac{x^2}{e^x} = 0$.
        De plus, $\lim_{x \to +\infty} \dfrac{2}{e^x} = 0$.
        Le facteur entre parenthèses tend donc vers : $0 + 0 - 1 = -1$.
        \item Comme $\lim_{x \to +\infty} e^x = +\infty$, on conclut par produit de limites :
        \[
        \boxed{\lim_{x \to +\infty} h_1(x) = -\infty}.
        \]
    \end{itemize}

    \item \textbf{Limite de $h_2(x) = \dfrac{2e^x - x}{x^2}$ en $+\infty$ :}
    \begin{itemize}
        \item Forme indéterminée « $\frac{\infty}{\infty}$ ».
        \item Scindons la fraction en deux termes pour tout $x > 0$ :
        \[
        h_2(x) = \frac{2e^x}{x^2} - \frac{x}{x^2} = 2 \frac{e^x}{x^2} - \frac{1}{x}.
        \]
        \item Par le théorème des croissances comparées :
        \[
        \lim_{x \to +\infty} \frac{e^x}{x^2} = +\infty \implies \lim_{x \to +\infty} 2\frac{e^x}{x^2} = +\infty.
        \]
        Comme $\lim_{x \to +\infty} \dfrac{1}{x} = 0$, on obtient par somme de limites :
        \[
        \boxed{\lim_{x \to +\infty} h_2(x) = +\infty}.
        \]
    \end{itemize}

    \item \textbf{Limite de $h_3(x) = \dfrac{e^x}{x^2 + 1}$ en $+\infty$ :}
    \begin{itemize}
        \item Forme indéterminée « $\frac{\infty}{\infty}$ ».
        \item Factorisons le dénominateur par son terme de plus haut degré $x^2$ :
        \[
        x^2 + 1 = x^2 \left(1 + \frac{1}{x^2}\right).
        \]
        Ainsi :
        \[
        h_3(x) = \frac{e^x}{x^2} \times \frac{1}{1 + \frac{1}{x^2}}.
        \]
        \item D'une part, par croissances comparées : $\lim_{x \to +\infty} \dfrac{e^x}{x^2} = +\infty$.
        D'autre part, $\lim_{x \to +\infty} \dfrac{1}{x^2} = 0$, donc $\lim_{x \to +\infty} \dfrac{1}{1 + \frac{1}{x^2}} = 1$.
        \item Par produit de limites :
        \[
        \boxed{\lim_{x \to +\infty} h_3(x) = +\infty}.
        \]
    \end{itemize}

    \item \textbf{Limite de $h_4(x) = e^{2x} - (x+1)e^x$ en $+\infty$ :}
    \begin{itemize}
        \item Forme indéterminée « $\infty - \infty$ ».
        \item Le terme dominant est $e^{2x} = (e^x)^2$, d'ordre d'exponentielle supérieur à $e^x$. Factorisons par $e^{2x}$ :
        \[
        h_4(x) = e^{2x} \left(1 - \frac{(x+1)e^x}{e^{2x}}\right) = e^{2x} \left(1 - \frac{x+1}{e^x}\right) = e^{2x} \left(1 - \frac{x}{e^x} - \frac{1}{e^x}\right).
        \]
        \item Par croissances comparées : $\lim_{x \to +\infty} \dfrac{x}{e^x} = 0$, et $\lim_{x \to +\infty} \dfrac{1}{e^x} = 0$.
        Le terme entre parenthèses tend donc vers : $1 - 0 - 0 = 1$.
        \item Puisque $\lim_{x \to +\infty} e^{2x} = +\infty$, on a par produit :
        \[
        \boxed{\lim_{x \to +\infty} h_4(x) = +\infty}.
        \]
    \end{itemize}

    \item \textbf{Limite de $h_5(x) = \dfrac{\sqrt{e^x + 1}}{x}$ en $+\infty$ :}
    \begin{itemize}
        \item Forme indéterminée « $\frac{\infty}{\infty}$ ».
        \item Factorisons par $e^x$ sous le radical pour $x > 0$ :
        \[
        \sqrt{e^x + 1} = \sqrt{e^x \left(1 + e^{-x}\right)} = \sqrt{e^x} \times \sqrt{1 + e^{-x}} = e^{x/2} \sqrt{1 + e^{-x}}.
        \]
        L'expression devient :
        \[
        h_5(x) = \frac{e^{x/2}}{x} \times \sqrt{1 + e^{-x}}.
        \]
        \item Pour évaluer $\dfrac{e^{x/2}}{x}$, effectuons le changement de variable :
        \[
        u = \frac{x}{2} \iff x = 2u.
        \]
        Lorsque $x \to +\infty$, $u \to +\infty$. Il s'ensuit que :
        \[
        \frac{e^{x/2}}{x} = \frac{e^u}{2u} = \frac{1}{2} \times \frac{e^u}{u}.
        \]
        D'après les croissances comparées, $\lim_{u \to +\infty} \dfrac{e^u}{u} = +\infty$, donc :
        \[
        \lim_{x \to +\infty} \frac{e^{x/2}}{x} = +\infty.
        \]
        \item Comme $\lim_{x \to +\infty} e^{-x} = 0$, on a $\lim_{x \to +\infty} \sqrt{1 + e^{-x}} = 1$.
        Par produit :
        \[
        \boxed{\lim_{x \to +\infty} h_5(x) = +\infty}.
        \]
    \end{itemize}

    \item \textbf{Limite de $h_6(x) = \dfrac{e^{2x} + x^2}{x^2 + x - 3}$ en $+\infty$ :}
    \begin{itemize}
        \item Forme indéterminée « $\frac{\infty}{\infty}$ ».
        \item Factorisons par le terme prépondérant au numérateur ($e^{2x}$) et au dénominateur ($x^2$) pour tout $x > 0$ :
        \[
        h_6(x) = \frac{e^{2x}\left(1 + \frac{x^2}{e^{2x}}\right)}{x^2\left(1 + \frac{1}{x} - \frac{3}{x^2}\right)} = \frac{e^{2x}}{x^2} \times \frac{1 + \frac{x^2}{e^{2x}}}{1 + \frac{1}{x} - \frac{3}{x^2}}.
        \]
        \item Étudions le comportement de chaque facteur :
        \begin{itemize}
            \item Pour $\dfrac{e^{2x}}{x^2}$ : Posons $u = 2x \iff x = \dfrac{u}{2}$, d'où $x^2 = \dfrac{u^2}{4}$.
            \[
            \frac{e^{2x}}{x^2} = \frac{e^u}{\frac{u^2}{4}} = 4 \frac{e^u}{u^2}.
            \]
            Par croissances comparées, $\lim_{u \to +\infty} \dfrac{e^u}{u^2} = +\infty$, donc $\lim_{x \to +\infty} \dfrac{e^{2x}}{x^2} = +\infty$.
            \item Par passage à l'inverse : $\lim_{x \to +\infty} \dfrac{x^2}{e^{2x}} = \lim_{u \to +\infty} \dfrac{u^2}{4e^u} = 0$.
            \item Pour le second facteur rationnel :
            \[
            \lim_{x \to +\infty} \frac{1 + \frac{x^2}{e^{2x}}}{1 + \frac{1}{x} - \frac{3}{x^2}} = \frac{1 + 0}{1 + 0 - 0} = 1.
            \]
        \end{itemize}
        \item Par produit de limites :
        \[
        \boxed{\lim_{x \to +\infty} h_6(x) = +\infty}.
        \]
    \end{itemize}
\end{enumerate}
\end{correctionbox}

% =========================================================================
% EXERCICE 12
% =========================================================================
\subsection*{Exercice 12 : Démonstration constructive des croissances comparées}
\addcontentsline{toc}{subsection}{Exercice 12 : Démonstration constructive des croissances comparées}

Le but de cet exercice est d'établir de manière autonome et rigoureuse le théorème des croissances comparées sans admettre le résultat.
\begin{enumerate}[label=\textbf{\arabic*.}]
    \item Démontrer que pour tout réel $x \in \R$, $e^x > x$.
    \item On considère la fonction $f$ définie sur $\R$ par :
    \[
    f(x) = e^x - \frac{x^2}{2}.
    \]
    Justifier la dérivabilité de $f$ sur $\R$, calculer $f'(x)$ et démontrer que $f$ est \textbf{strictement croissante} sur $\R$.
    \item En déduire que pour tout réel $x > 0$, on a :
    \[
    e^x > \frac{x^2}{2}.
    \]
    \item En déduire par comparaison la limite fondamentale :
    \[
    \lim_{x \to +\infty} \frac{e^x}{x} = +\infty.
    \]
    \item \textbf{Prolongement CPGE :} Expliquer comment cette méthode constructive se généralise pour établir que $\lim_{x \to +\infty} \dfrac{e^x}{x^n} = +\infty$ pour tout entier $n \ge 1$.
\end{enumerate}

\begin{correctionbox}[Exercice 12]
\begin{enumerate}[label=\textbf{\arabic*.}]
    \item \textbf{Démonstration de $e^x > x$ pour tout $x \in \R$ :}
    Considérons la fonction d'écart $\psi : \R \to \R$ définie par :
    \[
    \psi(x) = e^x - x - 1.
    \]
    $\psi$ est dérivable sur $\R$ et $\psi'(x) = e^x - 1$.
    \begin{itemize}
        \item Pour $x < 0$, $e^x < 1 \implies \psi'(x) < 0$.
        \item Pour $x > 0$, $e^x > 1 \implies \psi'(x) > 0$.
    \end{itemize}
    La fonction $\psi$ est donc strictement décroissante sur $]-\infty, 0]$ et strictement croissante sur $[0, +\infty[$.
    Elle admet en $x = 0$ un minimum absolu qui vaut :
    \[
    \psi(0) = e^0 - 0 - 1 = 0.
    \]
    Par suite, pour tout $x \in \R$, $\psi(x) \ge 0$, c'est-à-dire :
    \[
    e^x \ge x + 1.
    \]
    Comme pour tout réel $x$, $x + 1 > x$, on en déduit par stricte transitivité :
    \[
    \boxed{\forall x \in \R, \quad e^x > x}.
    \]

    \item \textbf{Étude de $f(x) = e^x - \dfrac{x^2}{2}$ et stricte croissance :}
    \begin{itemize}
        \item $f$ est la différence de la fonction exponentielle et d'un monôme du second degré, toutes deux dérivables sur $\R$. $f$ est donc dérivable sur $\R$.
        \item Calculons la dérivée de $f$ pour tout réel $x$ :
        \[
        f'(x) = e^x - \frac{2x}{2} = e^x - x.
        \]
        \item D'après la question 1, pour tout $x \in \R$, on a $e^x > x$, ce qui équivaut à :
        \[
        e^x - x > 0.
        \]
        Par conséquent :
        \[
        \forall x \in \R, \quad f'(x) > 0.
        \]
        La dérivée étant strictement positive sur l'intervalle $\R$, la fonction $f$ est \textbf{strictement croissante sur $\R$}.
    \end{itemize}

    \item \textbf{Minoration de $e^x$ pour $x > 0$ :}
    Puisque $f$ est strictement croissante sur $\R$, pour tout réel $x > 0$, on a l'inégalité stricte :
    \[
    f(x) > f(0).
    \]
    Calculons $f(0)$ :
    \[
    f(0) = e^0 - \frac{0^2}{2} = 1 - 0 = 1.
    \]
    Ainsi, pour tout $x > 0$ :
    \[
    e^x - \frac{x^2}{2} > 1 > 0.
    \]
    En ajoutant $\dfrac{x^2}{2}$ aux deux membres, on obtient immédiatement :
    \[
    \boxed{\forall x > 0, \quad e^x > \frac{x^2}{2}}.
    \]

    \item \textbf{Preuve par comparaison de $\lim_{x \to +\infty} \dfrac{e^x}{x} = +\infty$ :}
    Soit $x > 0$. Divisons l'inégalité précédente par le réel strictement positif $x$ (l'ordre est conservé) :
    \[
    \frac{e^x}{x} > \frac{\frac{x^2}{2}}{x} = \frac{x}{2}.
    \]
    Or :
    \[
    \lim_{x \to +\infty} \frac{x}{2} = +\infty.
    \]
    D'après le \textbf{théorème de comparaison par minoration} (théorème des gendarmes unilatéral), il vient directement :
    \[
    \boxed{\lim_{x \to +\infty} \frac{e^x}{x} = +\infty}.
    \]
    Le théorème fondamental des croissances comparées est ainsi rigoureusement démontré de manière constructive !

    \item \textbf{Prolongement CPGE : Généralisation inductive et formule de Taylor :}
    Cette démarche s'étend de manière naturelle par récurrence sur l'exposant polynomial.
    Pour tout entier $n \ge 1$, définissons la fonction d'écart à l'ordre $n$ :
    \[
    f_n(x) = e^x - \sum_{k=0}^{n} \frac{x^k}{k!} = e^x - \left(1 + x + \frac{x^2}{2!} + \dots + \frac{x^n}{n!}\right).
    \]
    Remarquons que pour tout $n \ge 1$, $f_n'(x) = f_{n-1}(x)$ et $f_n(0) = 0$.
    Par une récurrence immédiate :
    \begin{itemize}
        \item Puisque $f_1(x) = e^x - (1+x) > 0$ pour tout $x > 0$,
        \item $f_2'(x) = f_1(x) > 0 \implies f_2(x) > f_2(0) = 0 \implies e^x > 1 + x + \frac{x^2}{2} > \frac{x^2}{2}$,
        \item Par induction, pour tout entier $n \ge 1$ et pour tout $x > 0$ :
        \[
        e^x > \frac{x^{n+1}}{(n+1)!}.
        \]
    \end{itemize}
    En divisant par $x^n > 0$, il vient pour tout $x > 0$ :
    \[
    \frac{e^x}{x^n} > \frac{x}{(n+1)!} \xrightarrow[x \to +\infty]{} +\infty.
    \]
    On établit ainsi avec la plus grande élégance que l'exponentielle l'emporte sur n'importe quel monôme de degré arbitrairement grand :
    \[
    \forall n \in \N^*, \quad \lim_{x \to +\infty} \frac{e^x}{x^n} = +\infty \qquad \text{et} \qquad \lim_{x \to -\infty} x^n e^x = 0.
    \]
    En classe préparatoire MPSI, ce résultat s'interprète directement via la formule de Taylor avec reste intégral :
    \[
    e^x = \sum_{k=0}^n \frac{x^k}{k!} + \int_0^x \frac{(x-t)^n}{n!} e^t \, dt,
    \]
    dont l'intégrale positive pour $x > 0$ fournit la minoration polynomiale recherchée.
\end{enumerate}
\end{correctionbox}

\vspace{0.8cm}
\hrule
\vspace{0.8cm}
\end{document}
